\documentclass[12pt]{article}

\usepackage[a4paper, margin=1in]{geometry}
\usepackage{amsmath,amssymb,amsthm}
\usepackage{mathtools}
\usepackage{mathrsfs}
\usepackage{enumitem}
\usepackage{xcolor}
\usepackage[hidelinks]{hyperref}
\usepackage{fontspec}

\defaultfontfeatures{Ligatures=TeX}
\IfFontExistsTF{Times New Roman}
  {\setmainfont{Times New Roman}}
  {\setmainfont{TeX Gyre Termes}}
\IfFontExistsTF{Menlo}
  {\setmonofont{Menlo}}
  {\setmonofont{Latin Modern Mono}}

\newtheorem{theorem}{Theorem}[section]
\newtheorem{proposition}[theorem]{Proposition}
\newtheorem{lemma}[theorem]{Lemma}
\newtheorem{corollary}[theorem]{Corollary}
\theoremstyle{definition}
\newtheorem{definition}[theorem]{Definition}
\theoremstyle{remark}
\newtheorem{remark}[theorem]{Remark}

\newcommand{\tightlist}{%
  \setlength{\itemsep}{0pt}\setlength{\parskip}{0pt}}

\newcommand{\FISourceContinuityPhrase}{pointwise \(\sigma\)-weakly continuous}
\newcommand{\FIAppendixContinuityPhrase}{pointwise \(\sigma\)-weakly continuous}
\newcommand{\FIHeisenbergContinuityRemark}{%
\begin{remark}[Heisenberg-picture continuity convention]
The semigroup \(\{\mathcal E_t\}\) acts on observables, so the
continuity assumption in Definition~\ref{sec:ledger} is pointwise
\(\sigma\)-weak continuity on \(\mathcal A_U\), not norm continuity on
the ambient algebra. This is the natural continuity regime for the
Heisenberg-picture QDS language used throughout the paper.
\end{remark}}
\newcommand{\FIDisplayDate}{April 2026\\ \small Version 2}

\title{%
  \textbf{Foundation I: From Source to Manifestation}\\
  \textbf{A Conditional Same-Root Program for Geometry and Modular Flow}%
}

\author{
  Sidong Liu, PhD \\
  \small iBioStratix Ltd \\
  \small \texttt{sidongliu@hotmail.com}
}

\date{\FIDisplayDate}

\begin{document}
\emergencystretch=2em
\raggedbottom
\hbadness=5000
\vbadness=5000

\maketitle

\begin{abstract}
Paper~I of the Consciousness Trilogy~\cite{PaperI} begins from an
admissible accessible algebra \((\mathcal A_c,\omega)\) and derives the
inseparability of geometry and modular flow. The deeper source-side
question is how one reaches such an object from the pre-observer
substrate developed across Genesis, especially the HAFF and T-DOME
papers. The present paper, Foundation~I of the Consciousness Series,
does not claim a completed unified-source theorem. Instead, it
formulates an honest source-side foundation program and pushes its
currently closed parts into paper form.

We first state the correct conditional target: if a persistent
Nakajima--Zwanzig projection exists in the source datum, and if the
foreground-to-algebra transfer can be carried out, then the chain
\[
S_{\mathrm{source}} \to \mathcal K_* \to \mathfrak F \to (\mathcal A_c,\omega)
\to (\mathrm{Geom},\mathrm{Mod})
\]
is forced. Within that program, we prove the packaging result that
budget violation in the persistent record process forces
symmetry-broken foreground compression, thereby rewriting the T-DOME
I--II chain as a clean Foundation~I proposition. We then prove a first
restricted PS2 closure in a product-classical, decoherence-dominated
pointer regime: pointer coarse-graining yields fragment-local omission
losses, semigroup decay yields a uniformly bounded coherence remainder,
and a weight gap larger than the resulting remainder forces exact
top-\(k^*\) pointer selection. In a first finite-dimensional faithful
pointer-support regime, the required coherence constants become
explicit in Hilbert--Schmidt norm. We then define a
pointer-aligned realization of the foreground frame inside the ambient
algebra, construct the associated pointer-generated von Neumann
algebra, and replace the earlier open-ended transfer candidate by a
minimal modular-decoherence stable hull
\(\Phi_{\mathfrak F}^{\mathrm{st}}(\mathfrak F)\). This stable-hull
transfer is well-defined, yields a von Neumann subalgebra, and carries a
faithful restricted state.

Finally, under the exact fixed-point form of the HAFF accessibility
condition, we prove the inclusion-level comparison
\[
\Phi_{\mathfrak F}^{\mathrm{st}}(\mathfrak F)\subseteq
\mathcal A_c^{\mathrm{HAFF}},
\]
while leaving equality as an open strong-form problem. The upshot is
deliberately calibrated: Genesis already supplies substantial antecedent
support for Foundation~I, but the completed deductive source-side
foundation still depends on two genuinely open interfaces, namely
persistent projection selection and the model-independent derivation of
PS2 beyond the restricted product-classical pointer regime treated here.
\end{abstract}

\clearpage
\tableofcontents
\clearpage

\section{Introduction}

Paper~I of the Consciousness Trilogy~\cite{PaperI} starts at the
manifest layer: one assumes an admissible accessible algebra
\((\mathcal A_c,\omega)\), proves that geometry and modular flow are not
independently deformable, and then studies the consequences of that
inseparability. That was the correct place for the first trilogy paper
to begin, but it leaves a source-foundation question unresolved. What, in
the broader Genesis programme, plays the role of the source-side
antecedent to \((\mathcal A_c,\omega)\)?

The present paper, Foundation~I of the Consciousness Series, addresses
exactly that question. It does not claim that Genesis already contains a
completed theorem deriving the trilogy from a single source object. What
it does claim is narrower and more precise: the relevant source-side
foundation can now be decomposed cleanly, its nearly closed parts can be
written as formal propositions, and its actual bottlenecks can be
localized without conflating vision and theorem.

The philosophical motivation is correspondingly limited. Paper~I showed
that geometry and modular flow are not independently deformable once one
is already at the manifest pair \((\mathcal A_c,\omega)\).
Foundation~I asks the prior structural question: why should those two
readings belong together at all? In the present paper, the same-root
idea means only that one seeks a source-side chain from which geometric
and modular aspects later arise together. This paper treats that
mathematical bridge from source data to accessible algebra; it does not
attempt to settle consciousness ontology or the broader philosophy of
mind.

The correct three-layer structure is
\[
S_{\mathrm{source}} \longrightarrow (\mathcal A_c,\omega)
\longrightarrow (\mathrm{Geom},\mathrm{Mod}),
\]
with an important internal refinement:
\[
S_{\mathrm{source}}
\longrightarrow
\mathcal O_A
\longrightarrow
\mathcal O_B
\longrightarrow
(\mathcal A_c,\omega).
\]
Here \(S_{\mathrm{source}}\) is the pre-observer source datum,
\(\mathcal O_A\) is a persistent projected record process with
non-Markovian memory, and \(\mathcal O_B\) is the symmetry-broken
foreground frame produced when full-rank persistence exceeds the
available processing budget.

Four claims structure the paper.
\begin{enumerate}
\item We state the honest \emph{conditional} Foundation~I target. The
  two genuine open points are persistent projection selection and the
  model-independent foreground-alignment step PS2 that makes the
  transfer construction canonical; the existence of an exact fixed-point
  pointer sector is treated only as an exact-regime PS1 idealization.
\item We promote the T-DOME I--II rate chain to a formal proposition:
  budget violation forces symmetry-broken foreground compression.
\item We prove a first restricted PS2 closure package: in a
  product-classical, decoherence-dominated pointer regime, pointer
  coarse-graining yields fragment-local omission losses, semigroup decay
  yields a uniformly bounded coherence remainder, and a weight gap
  larger than the resulting remainder forces exact top-\(k^*\) pointer
  selection.
\item We give the paper's main new operator-algebraic construction: a
  stable-hull realization of the foreground sector. Concretely, we
  construct a minimal modular-decoherence stable hull, prove that it is
  a von Neumann subalgebra with faithful restricted state, and prove the
  first comparison theorem
  \(\mathcal A_{\mathfrak F}\subseteq
  \mathcal A_c^{\mathrm{HAFF}}\) under the exact fixed-point
  accessibility condition.
\end{enumerate}

This is enough to fix the current calibration. Genesis provides
substantial antecedent support for Foundation~I. It does not yet
provide a finished deductive source-side foundation. Equality between the
foreground-generated algebra and the full HAFF accessible algebra is not
claimed here and should remain an open strong-form problem.

The paper is organized as follows. Section~\ref{sec:ledger} defines the
source datum and the stage outputs. Section~\ref{sec:target} states the
target conditional theorem and isolates the source-selection problem.
Section~\ref{sec:budget} proves the budget-forced compression
proposition. Section~\ref{sec:foreground} defines pointer alignment,
foreground realization, and the PS2 target. Section~\ref{sec:ps2model}
proves a first restricted PS2 closure package in the
product-classical, decoherence-dominated pointer regime, and records a
bookkeeping lemma fixing the diagonal/off-sector split needed for the
later model-independent PS2 programme.
Section~\ref{sec:haff} proves the first inclusion-level comparison with
HAFF accessibility and states a partial Lemma~C.
Section~\ref{sec:open} records the remaining open problems.
Appendix~\ref{app:premises} summarizes the external Genesis inputs used
in the paper in self-contained mathematical form, and
Appendix~\ref{app:conditional_ledger} collects the
named conditional interfaces that remain open.

\section{Source Datum and Object Ledger}\label{sec:ledger}

\begin{definition}[Source datum]
A \emph{source datum} is a quadruple
\[
S_{\mathrm{source}}
:=
(\mathcal H_U,\mathcal A_U,\Omega,\{\mathcal E_t\}_{t\ge 0}),
\]
where:
\begin{enumerate}
\item \((\mathcal H_U,\mathcal A_U,\Omega)\) is a Type~III\(_1\) von
  Neumann algebra in standard form, with \(\Omega\) cyclic and
  separating for \(\mathcal A_U\);
\item \(\{\mathcal E_t\}_{t\ge 0}\) is a \FISourceContinuityPhrase\
  semigroup of normal completely positive unital maps on
  \(\mathcal A_U\).
\end{enumerate}
\end{definition}

\FIHeisenbergContinuityRemark

The ambient standard-form side is the HAFF input
\cite{HAFFD}; the projected open-system side is expressed in
Nakajima--Zwanzig language \cite{Nakajima1958,Zwanzig1960,
BreuerPetruccione2002}.

\begin{definition}[Stage outputs]
Given a source datum \(S_{\mathrm{source}}\), we use the following
ledger.
\begin{enumerate}
\item \textbf{Stage A output}
  \[
  \mathcal O_A
  :=
  (\mathcal P,\omega_{\mathcal P,t},\mathcal K_*(t,s)),
  \]
  where \(\mathcal P:\mathcal A_U\to\mathcal A_U\) is a
  Nakajima--Zwanzig projection,
  \(\omega_{\mathcal P,t}:=\mathcal P_*(\omega_t)\) is the projected
  normal state on the predual, and \(\mathcal K_*(t,s)\) is the dual
  memory kernel.
\item \textbf{Stage B output}
  \[
  \mathcal O_B
  :=
  (\mathcal P,\omega_{\mathcal P,t},\mathcal K_*(t,s),\mathfrak F,\pi_{\mathrm{ego}}),
  \]
  where \(\mathfrak F\) is a \(k^*\)-dimensional foreground frame and
  \(\pi_{\mathrm{ego}}\) is the ego prior carried forward from the
  broken phase of T-DOME II and into the self-calibration analysis of
  T-DOME III~\cite{TDOMEII,TDOMEIII}.
\item \textbf{Stage C output}
  \[
  \mathcal O_C
  :=
  (\mathcal A_c,\omega),
  \qquad
  \omega=\langle\Omega,\cdot\,\Omega\rangle|_{\mathcal A_c}.
  \]
\end{enumerate}
\end{definition}

\begin{remark}[Stage D]
Once an admissible \((\mathcal A_c,\omega)\) is available, the
Paper~I chain from accessible algebra to jointly constrained geometry
and modular flow is already in place~\cite{PaperI}. The present note is
therefore concerned only with Stages A--C.
\end{remark}

\section{Target Conditional Theorem and the Source-Selection Problem}\label{sec:target}

The honest Foundation~I target is conditional and should be stated that
way from the start.

\paragraph{Target conditional theorem.}
Let
\[
S_{\mathrm{source}}
=
(\mathcal H_U,\mathcal A_U,\Omega,\{\mathcal E_t\}_{t\ge 0})
\]
be a source datum. Assume:
\begin{enumerate}
\item there exists a persistent Nakajima--Zwanzig projection
  \(\mathcal P\) whose projected predual dynamics carries a
  non-Markovian kernel \(\mathcal K_*(t,s)\);
\item the induced record process exceeds its full-rank processing budget
  and therefore forces a compressed foreground frame \(\mathfrak F\);
\item there exists a foreground embedding
  \[
  \iota_{\mathfrak F}:V_{\mathrm{fg}}(\mathfrak F)\hookrightarrow
  \mathcal A_U^{\mathrm{sa}}
  \]
  from which the transfer construction can be carried out and shown to
  agree, at least at the inclusion level, with the HAFF accessibility
  construction;
\item the resulting \((\mathcal A_c,\omega)\) satisfies the
  admissibility hypotheses needed by Paper~I.
\end{enumerate}
Then the chain
\[
S_{\mathrm{source}}
\to
\mathcal O_A
\to
\mathcal O_B
\to
(\mathcal A_c,\omega)
\to
(\mathrm{Geom},\mathrm{Mod})
\]
is forced.

\begin{remark}[What is actually open]
The conditional theorem is open in exactly two substantive places. The
first is the existence and selection of a persistent projection
\(\mathcal P\). The second is the full transfer problem from the
foreground frame to the operator-algebraic accessibility construction.
The first is an observer-emergence problem. The second is the main new
interface problem.
\end{remark}

\begin{lemma}[Conditional source-to-kernel interface]\label{lem:A}
Let \(S_{\mathrm{source}}\) be a source datum. Suppose there exists a
Nakajima--Zwanzig projection \(\mathcal P\) such that the projected
predual dynamics admits a non-Markovian kernel \(\mathcal K_*(t,s)\)
and yields a persistent record process. Then
\[
\mathcal O_A=(\mathcal P,\omega_{\mathcal P,t},\mathcal K_*(t,s))
\]
satisfies the input conditions of Stage~B.
\end{lemma}

\begin{proof}
This is the conditional form of the source-side interface. The assumed
projection already provides the Stage~A output. Within the T-DOME~I
modeling framework~\cite{TDOMEI}, the persistence assumption is analyzed
in a non-Markovian Nakajima--Zwanzig regime; the present paper imports
that framing rather than claiming it as a general open-quantum-systems
theorem. The point is not to solve the selection problem
constructively, but to separate that problem from the downstream
thermodynamic and algebraic consequences.
\end{proof}

\begin{remark}[Constructive version remains open]
Lemma~\ref{lem:A} is intentionally conditional. The missing constructive
question is why a given source datum selects a persistent projection at
all. Nothing in the present note should be read as claiming that this
observer-emergence problem has been solved.
\end{remark}

\section{Budget-Forced Foreground Compression}\label{sec:budget}

We now formalize the part of Foundation~I that is essentially
already present across T-DOME I and II, but not yet stated in the
present source-to-accessibility language.

\begin{proposition}[Budget-forced foreground compression]\label{prop:budget}
Let
\[
\mathcal O_A=(\mathcal P,\omega_{\mathcal P,t},\mathcal K_*(t,s))
\]
be a persistent record process arising from a source datum
\(S_{\mathrm{source}}\). Write
\[
D:=\operatorname{rank}(\mathcal P)
\]
for the effective number of active record components, and let
\(h_\mu>0\) be the entropy rate per component. Assume:
\begin{enumerate}
\item the induced record process satisfies the T-DOME II hypotheses
  corresponding to the Computational Ceiling and Compression Theorem
  \cite{TDOMEII};
\item the ceiling binds in the sense that the symmetric full-rank
  processing requirement exceeds the available budget,
  \[
  h_\mu D>\dot{\mathcal C}_{\mathrm{budget}};
  \]
\item the survival-distortion landscape is non-degenerate in the sense
  required by the T-DOME II symmetry-breaking theorem.
\end{enumerate}
Then every survival-optimal strategy is realized by a foreground frame
\(\mathfrak F\) with
\[
k^*
=
\Big\lfloor
\frac{\dot{\mathcal C}_{\mathrm{budget}}}{h_\mu}
\Big\rfloor
<D,
\]
and the passage from full-rank tracking to the
\(k^*\)-component foreground/background split is spontaneous symmetry
breaking.
\end{proposition}

\begin{proof}
By the Computational Ceiling of T-DOME~II~\cite{TDOMEII}, any symmetric
strategy that tracks all \(D\) active components requires processing
rate at least \(h_\mu D\). Under assumption~(2), this exceeds the
available budget, so the full-rank symmetric strategy is infeasible. The
Compression Theorem of T-DOME~II then identifies the survival-optimal
admissible rank as
\[
k^*
=
\Big\lfloor
\frac{\dot{\mathcal C}_{\mathrm{budget}}}{h_\mu}
\Big\rfloor,
\]
with the optimum saturating the budget. Under assumption~(3), the
Necessity of Symmetry Breaking theorem upgrades this forced compression
into the selection of a privileged frame \(\mathfrak F\) and a
corresponding foreground/background decomposition. Since \(k^*<D\), the
resulting representation is genuinely compressed.

The role of T-DOME~I~\cite{TDOMEI} is to supply the thermodynamic reason
that persistent non-Markovian processes are driven into the budget
binding regime in the first place: positive entropy production and
finite processing resources make indefinite full-rank retention
unsustainable. Hence budget violation forces symmetry-broken foreground
compression.
\end{proof}

\begin{remark}[Status of Proposition~\ref{prop:budget}]
The proposition is best read as a notation-translation theorem. Its
substance is already contained in the T-DOME I--II chain; what was
missing was a clean statement in the present Foundation~I language.
\end{remark}

\section{Foreground Realization inside the Ambient Algebra}\label{sec:foreground}

Lemma~C requires a genuinely new object: a way to realize the
foreground frame inside the ambient algebra.

\begin{definition}[Pointer-aligned observable]\label{def:pointer}
A self-adjoint observable \(O\in\mathcal A_U^{\mathrm{sa}}\) is called
\emph{pointer-aligned} with respect to the decoherence semigroup
\(\{\mathcal E_t\}_{t\ge 0}\) if it is asymptotically fixed:
\[
\lim_{t\to\infty}\|\mathcal E_t(O)-O\|=0.
\]
In the strong form one requires exact invariance,
\[
\mathcal E_t(O)=O
\qquad
\forall\, t\ge 0.
\]
\end{definition}

\begin{remark}
The operator-norm asymptotic notion is the working definition used in
this paper's exact dephasing-dominated regime. It is not claimed to be
the most natural decoherence notion in general: for a generic
Heisenberg-picture semigroup, weak*, strong, or predual-level
convergence criteria may be more appropriate. The exact version is used
when comparing with the HAFF fixed-point algebra. This is also the point
where Foundation~I touches the pointer-state and quantum Darwinism
literature \cite{Zurek2003,Zurek2009}; the exact fixed-point condition
should be understood as an idealization of the approximate einselection
described there.
\end{remark}

Let
\[
V_{\mathrm{fg}}(\mathfrak F)=\operatorname{span}\{f_1,\dots,f_{k^*}\}
\]
be an orthonormal foreground basis selected by the T-DOME II
rate-distortion optimum. Guided by the pointer basis selected in the
broken phase \cite{TDOMEII}, choose self-adjoint ambient observables
\[
O_1,\dots,O_{k^*}\in\mathcal A_U^{\mathrm{sa}}
\]
and define
\[
\iota_{\mathfrak F}(f_i):=O_i,
\qquad
\iota_{\mathfrak F}\Big(\sum_{i=1}^{k^*} a_i f_i\Big)
:=
\sum_{i=1}^{k^*} a_i O_i.
\]
Write
\[
\mathfrak S_{\mathfrak F}
:=
\operatorname{span}\{1,O_1,\dots,O_{k^*}\}
\subset\mathcal A_U
\]
for the resulting operator system and
\[
\mathcal A_{\mathfrak F}:=W^*(\mathfrak S_{\mathfrak F})
\]
for its von Neumann envelope.

\begin{proposition}[Pointer realization of the foreground sector]\label{prop:pointer}
Let \(\mathfrak F\) be a T-DOME foreground frame with
\[
V_{\mathrm{fg}}(\mathfrak F)
=
\operatorname{span}\{f_1,\dots,f_{k^*}\}.
\]
Assume there exists a family of self-adjoint ambient observables
\[
O_1,\dots,O_{k^*}\in\mathcal A_U^{\mathrm{sa}}
\]
such that:
\begin{enumerate}
\item \(O_1,\dots,O_{k^*}\) are linearly independent;
\item each \(O_i\) is pointer-aligned under \(\{\mathcal E_t\}\);
\item the ordering of their survival relevance matches the ordering
  selected by the T-DOME foreground optimization.
\end{enumerate}
Then the linear extension \(f_i\mapsto O_i\) defines an injective map
\[
\iota_{\mathfrak F}:V_{\mathrm{fg}}(\mathfrak F)\hookrightarrow
\mathcal A_U^{\mathrm{sa}}.
\]
Moreover,
\[
\mathfrak S_{\mathfrak F}
=
\operatorname{span}\{1,O_1,\dots,O_{k^*}\}
\]
is a finite-dimensional operator system, and
\[
\mathcal A_{\mathfrak F}
=
W^*(\mathfrak S_{\mathfrak F})
\]
is a von Neumann subalgebra of \(\mathcal A_U\).
\end{proposition}

\begin{proof}
Linear independence of \(O_1,\dots,O_{k^*}\) makes the linear extension
of \(f_i\mapsto O_i\) injective. Since the \(O_i\) are self-adjoint,
\(\mathfrak S_{\mathfrak F}\) is self-adjoint and contains the
identity, hence is an operator system. By definition,
\(W^*(\mathfrak S_{\mathfrak F})\) is the von Neumann algebra generated
by that operator system inside \(\mathcal A_U\). Therefore
\(\mathcal A_{\mathfrak F}\) is a von Neumann subalgebra of
\(\mathcal A_U\).
\end{proof}

The next refinement is to replace the bare existence assumption on the
ambient observables \(O_i\) by a more structured compatibility package
that points directly back to T-DOME I and II.

\paragraph{The PS1--PS3 package.}
The structured foreground-realization input naturally separates into
three pieces:
\begin{enumerate}
\item[(PS1)] in the exact dephasing-dominated regime, the decoherence
  semigroup admits a commuting exact fixed-point pointer sector inside
  \(\mathcal A_U\);
\item[(PS2)] the T-DOME sufficient-statistic coordinates agree with the
  coordinates induced by that pointer sector;
\item[(PS3)] the T-DOME rate-distortion optimum selects an ordered
  subset of pointer channels by survival relevance.
\end{enumerate}
In this decomposition, PS1 is an exact-regime idealization of the
decoherence-side input (not a consequence of generic decoherence or
quantum Darwinism alone, but of the noiseless-sector limit in which
pointer projections are exactly fixed rather than approximately
einselected), whereas PS2 is the genuine model-independent bridge
bottleneck and PS3 is downstream from it.

\begin{proposition}[Conditional pointer-sector realization]\label{prop:pointersector}
Assume:
\begin{enumerate}
\item[(PS1)] in the exact dephasing-dominated regime, the decoherence
  semigroup admits a commuting family of nonzero exact fixed-point
  projections
  \(\{P_1,\dots,P_D\}\subset \operatorname{Fix}(\mathcal E)\)
  generating a classical pointer sector
  \[
  \mathcal A_{\mathrm{ptr}}:=W^*(P_1,\dots,P_D)\subseteq \mathcal A_U;
  \]
\item[(PS2)] the sufficient-statistic record coordinates used in T-DOME II are
  the coordinates induced by this pointer sector through the T-DOME I
  coarse-graining map to the classical sector;
\item[(PS3)] the T-DOME II rate-distortion optimum selects an ordered subset
  \(I^*=\{i_1,\dots,i_{k^*}\}\subseteq \{1,\dots,D\}\) by survival
  relevance.
\end{enumerate}
Let
\[
V_{\mathrm{fg}}(\mathfrak F)=\operatorname{span}\{f_1,\dots,f_{k^*}\}
\]
be ordered compatibly with \(I^*\), and define
\[
\iota_{\mathfrak F}(f_j):=P_{i_j}.
\]
Then:
\begin{enumerate}
\item \(\iota_{\mathfrak F}:V_{\mathrm{fg}}(\mathfrak F)\hookrightarrow
  \mathcal A_U^{\mathrm{sa}}\) is injective;
\item the image lies in the exact decoherence fixed-point algebra;
\item the induced pointer-generated algebra satisfies
  \[
  \mathcal A_{\mathfrak F}\subseteq \mathcal A_{\mathrm{ptr}}.
  \]
\end{enumerate}
\end{proposition}

\begin{proof}
By assumption~(1), each \(P_i\) is self-adjoint and satisfies
\(\mathcal E_t(P_i)=P_i\) for all \(t\ge 0\). Distinct nonzero
orthogonal projections are linearly independent, so the assignment
\(f_j\mapsto P_{i_j}\) extends to an injective linear map
\(\iota_{\mathfrak F}\). Exact fixedness gives the second claim. Since
the \(P_{i_j}\) lie in the commuting pointer sector
\(\mathcal A_{\mathrm{ptr}}\), the von Neumann algebra generated by
\(1,P_{i_1},\dots,P_{i_{k^*}}\) is contained in
\(\mathcal A_{\mathrm{ptr}}\).
\end{proof}

\begin{remark}[Status of PS1]
PS1 is the least problematic part of the package. In decoherence theory
and quantum Darwinism, robust classical records are carried by a
commuting pointer sector selected by the environment
\cite{Zurek2003,Zurek2009}. In the exact fixed-point regime used here,
this is precisely the statement that the decoherence semigroup admits a
commuting family of exact fixed-point projections generating a classical
pointer algebra. HAFF Paper~D realizes this exactly in the
dephasing-dominated finite-dimensional model, where the diagonal
fixed-point algebra is the pointer basis algebra \cite{HAFFD}. The real
model-independent unresolved content is therefore PS2, with PS3
following once PS2 is established.
\end{remark}

\paragraph{The PS2 alignment problem.}
Proposition~\ref{prop:pointersector} sharpens the existence assumption
on \(\iota_{\mathfrak F}\), but it does not yet derive the compatibility
package from the broken-phase output itself in full generality.
Section~\ref{sec:ps2model} closes a first restricted version in a
product-classical pointer regime. The real unresolved content is
therefore the model-independent form of the middle condition: the claim
that the T-DOME sufficient-statistic coordinates and the pointer-sector
coordinates are the same optimization problem written in two languages.
That is the full PS2 problem.

\paragraph{Target lemma (PS2 alignment; open).}
Let
\[
\mathcal A_{\mathrm{ptr}}=W^*(P_1,\dots,P_D)
\]
be the pointer sector from Proposition~\ref{prop:pointersector}, and
define the pointer-coordinate process
\[
p_i(t):=\omega_t(P_i),\qquad i=1,\dots,D.
\]
For an ordered subset
\[
I=\{i_1,\dots,i_k\}\subseteq \{1,\dots,D\},
\]
write \(\mathcal F_I\) for the pointer-adapted frame retaining those
channels. A natural target theorem is the following.

Assume:
\begin{enumerate}
\item the T-DOME sufficient statistics \(\{c_i(t)\}_{i=1}^D\) agree,
  after reindexing, with the pointer-coordinate process
  \(\{p_i(t)\}_{i=1}^D\);
\item there exist nonnegative weights \(w_1,\dots,w_D\) and a constant
  \(D_0\) such that for every pointer-adapted frame \(\mathcal F_I\),
  \[
  D(\mathcal F_I)=D_0+\sum_{i\notin I} w_i;
  \]
\item admissible frames satisfy \(|I|=k^*=
  \lfloor \mathcal C_{\mathrm{budget}}/h_\mu\rfloor\).
\end{enumerate}
Then every survival-distortion minimizer retains exactly the
\(k^*\) channels of largest weight, up to ties. In particular, the
T-DOME foreground directions coincide with the top-ranked pointer
directions, and the ordered subset \(I^*\) of
Proposition~\ref{prop:pointersector} is selected canonically.

\paragraph{Proof strategy.}
This target should be attacked by an exchange argument. If an admissible
subset \(I\) omits some channel \(a\) and retains some channel \(b\)
with \(w_a>w_b\), then replacing \(b\) by \(a\) changes the distortion
by
\[
D(\mathcal F_{I\setminus\{b\}\cup\{a\}})-D(\mathcal F_I)
=-w_a+w_b<0,
\]
contradicting optimality. Once the distortion has been reduced to such a
channelwise omission law, alignment is immediate. The real analytic
burden is therefore to derive the weights \(w_i\) from the T-DOME
survival functional in a way compatible with the pointer-sector
coarse-graining.

\paragraph{First realistic route: approximate PS2b in the
decoherence-dominated regime.}
The exact separability law above is likely too strong as a first target
in the fully correlated setting. A more plausible theorem program is to
prove approximate separability once decoherence has suppressed the
cross-channel terms strongly enough that the pointer sector is
approximately classical.

Introduce a decoherence-depth parameter \(\tau_{\mathrm{dec}}\), and let
\(\varepsilon_{\mathrm{dec}}(\tau_{\mathrm{dec}})\to 0\) measure the
size of the residual cross-channel distortion after coarse-graining
beyond that scale. The natural asymptotic target is:
\[
D(\mathcal F_I)
=
D_0+\sum_{i\notin I} w_i + R_I,
\qquad
|R_I|\le \varepsilon_{\mathrm{dec}}(\tau_{\mathrm{dec}}).
\]
Under this bound, the same exchange argument gives
\[
D(\mathcal F_{I\setminus\{b\}\cup\{a\}})-D(\mathcal F_I)
\le
-(w_a-w_b)+2\varepsilon_{\mathrm{dec}}(\tau_{\mathrm{dec}}).
\]
Therefore any exchange across a weight gap larger than
\(2\varepsilon_{\mathrm{dec}}(\tau_{\mathrm{dec}})\) strictly improves
the objective. In particular, if the threshold gap
\[
\Delta_{k^*}:=w_{(k^*)}-w_{(k^*+1)}
\]
exceeds \(2\varepsilon_{\mathrm{dec}}(\tau_{\mathrm{dec}})\), then the
top-\(k^*\) pointer subset is forced, up to ties above the gap.

This is likely the right first Bridge~1 theorem target. It matches the
physics of decoherence, which suppresses off-diagonal couplings rather
than eliminating them exactly in full generality.
Section~\ref{sec:ps2model} closes this target in a first restricted
product-classical regime. The real analytic task is thus to prove that
the T-DOME survival distortion admits such a uniform remainder estimate
in the decoherence-dominated regime.

\paragraph{Analytic source of \(\varepsilon_{\mathrm{dec}}\).}
The remainder estimate cannot come from abstract operator algebra alone.
It has to be extracted from the concrete T-DOME I survival identity
\[
\beta\,\mathcal S
=
-\Delta I(S{:}E)
-\Delta D_{\mathrm{KL}}
-\beta\,\Delta\langle H_{\mathrm{ctrl}}\rangle.
\]
In the pointer basis, the natural next step is to decompose the record
process and the induced state into diagonal and off-diagonal parts,
\[
X(t)=X^{\mathrm{diag}}(t)+X^{\mathrm{off}}(t),
\qquad
\rho(t)=\rho_{\mathrm{diag}}(t)+\rho_{\mathrm{off}}(t),
\]
where the diagonal part is generated by the commuting pointer
projections and the off-diagonal part carries the residual coherences.

The target analytic lemma is then the following: in the
decoherence-dominated regime, the off-diagonal contribution to the
survival identity is uniformly bounded by a function
\(\varepsilon_{\mathrm{dec}}(\tau_{\mathrm{dec}})\to 0\), so that for
every pointer-adapted frame \(\mathcal F_I\),
\[
D(\mathcal F_I)
=
D_0+\sum_{i\notin I} w_i + R_I,
\qquad
|R_I|\le \varepsilon_{\mathrm{dec}}(\tau_{\mathrm{dec}}).
\]
The likely proof route is:
\begin{enumerate}
\item use the decoherence semigroup to obtain exponential or
  semigroup-gap suppression of the off-pointer sector;
\item propagate that suppression through the three terms in the survival
  identity;
\item identify the diagonal leading term as an additive pointer-channel
  omission loss.
\end{enumerate}
This identifies the analytic content that
Section~\ref{sec:ps2model} closes in a first restricted regime.

\section{A Restricted PS2 Result in the Decoherence-Dominated Pointer Regime}\label{sec:ps2model}

The full PS2 alignment problem remains open. Still, a genuine theorem
can already be proved in a restricted regime. The point of the present
section is not to claim a model-independent PS2 theorem, but to show
that in a product-classical, decoherence-dominated pointer regime, the
structural reduction to fragment-local omission losses and the estimate
reduction to a uniformly small coherence remainder can both be closed
explicitly, after which approximate separability and top-\(k^*\)
selection follow by a completely explicit argument.

\begin{proposition}[Restricted approximate PS2b in the product-classical pointer regime]\label{prop:ps2_restricted}
Fix \(k\in\{1,\dots,D-1\}\), and let \(\mathcal F_I\) range over all
pointer-adapted frames with \(|I|=k\). Assume the T-DOME survival
distortion uses the squared loss
\[
\ell(x)=x^2.
\]
Suppose that for each admissible \(I\) there exist random variables
\[
L_i=L_i(\xi),
\qquad
i=1,\dots,D,
\]
and a coherence remainder \(C_I=C_I(\xi)\) such that
\[
\mathcal S_{\mathrm{full}}(\xi)-\mathcal S_{\mathcal F_I}(\xi)
=
\sum_{i\notin I} L_i(\xi)+C_I(\xi)
\]
for almost every environmental realization \(\xi\). Assume further:
\begin{enumerate}
\item \textbf{Pairwise diagonal orthogonality:}
  \[
  \mathbb E_\xi[L_iL_j]=0
  \qquad
  \text{for } i\neq j;
  \]
\item \textbf{Uniform coherence bound:}
  \[
  \|C_I\|_{L^\infty(\xi)}
  \le
  \eta_{\mathrm{dec}}(\tau_{\mathrm{dec}})
  \qquad
  \text{for every admissible } I;
  \]
\item \textbf{Finite diagonal size:}
  \[
  \Lambda_k
  :=
  \sup_{|I|=k}
  \sum_{i\notin I}\|L_i\|_{L^2(\xi)}
  <
  \infty.
  \]
\end{enumerate}
Define
\[
w_i:=\|L_i\|_{L^2(\xi)}^2
\]
and
\[
\varepsilon_{\mathrm{dec}}(\tau_{\mathrm{dec}})
:=
2\Lambda_k\,\eta_{\mathrm{dec}}(\tau_{\mathrm{dec}})
+\eta_{\mathrm{dec}}(\tau_{\mathrm{dec}})^2.
\]
Then for every admissible \(I\),
\[
D(\mathcal F_I)
=
\sum_{i\notin I} w_i + R_I,
\qquad
|R_I|
\le
\varepsilon_{\mathrm{dec}}(\tau_{\mathrm{dec}}).
\]
Consequently, if
\[
\Delta_k:=w_{(k)}-w_{(k+1)}
>
2\varepsilon_{\mathrm{dec}}(\tau_{\mathrm{dec}}),
\]
then every minimizer of \(D(\mathcal F_I)\) retains exactly the
top-\(k\) pointer channels, up to ties above the gap.
\end{proposition}

\begin{proof}
By definition of the distortion and the squared loss,
\[
D(\mathcal F_I)
=
\mathbb E_\xi\!\left[
\left(
\sum_{i\notin I}L_i + C_I
\right)^2
\right].
\]
Expanding the square gives
\[
D(\mathcal F_I)
=
\sum_{i\notin I}\mathbb E_\xi[L_i^2]
+
2\sum_{\substack{i,j\notin I\\ i<j}}\mathbb E_\xi[L_iL_j]
+
2\sum_{i\notin I}\mathbb E_\xi[L_iC_I]
+
\mathbb E_\xi[C_I^2].
\]
The pairwise orthogonality assumption cancels the middle sum, so
\[
D(\mathcal F_I)
=
\sum_{i\notin I} w_i + R_I,
\]
where
\[
R_I
:=
2\sum_{i\notin I}\mathbb E_\xi[L_iC_I]
+
\mathbb E_\xi[C_I^2].
\]
Because the environmental expectation is taken against a probability
measure, \(\|C_I\|_{L^2(\xi)}\le \|C_I\|_{L^\infty(\xi)}\). Hence
\[
|\mathbb E_\xi[L_iC_I]|
\le
\|L_i\|_{L^2(\xi)}\|C_I\|_{L^2(\xi)}
\le
\|L_i\|_{L^2(\xi)}\eta_{\mathrm{dec}}(\tau_{\mathrm{dec}}),
\]
and therefore
\[
|R_I|
\le
2\eta_{\mathrm{dec}}(\tau_{\mathrm{dec}})
\sum_{i\notin I}\|L_i\|_{L^2(\xi)}
+
\eta_{\mathrm{dec}}(\tau_{\mathrm{dec}})^2
\le
\varepsilon_{\mathrm{dec}}(\tau_{\mathrm{dec}}).
\]
This proves the approximate separability law.

Now let \(I\) be admissible, and suppose \(a\notin I\), \(b\in I\) with
\(w_a>w_b\). Writing
\[
I':=I\setminus\{b\}\cup\{a\},
\]
the additive part changes by
\[
\sum_{i\notin I'}w_i-\sum_{i\notin I}w_i
=
-w_a+w_b.
\]
The remainder terms satisfy
\[
R_{I'}-R_I
\le
|R_{I'}|+|R_I|
\le
2\varepsilon_{\mathrm{dec}}(\tau_{\mathrm{dec}}),
\]
so
\[
D(\mathcal F_{I'})-D(\mathcal F_I)
\le
-(w_a-w_b)+2\varepsilon_{\mathrm{dec}}(\tau_{\mathrm{dec}}).
\]
Therefore any exchange across a weight gap strictly larger than
\(2\varepsilon_{\mathrm{dec}}(\tau_{\mathrm{dec}})\) lowers the
distortion. If
\(\Delta_k>2\varepsilon_{\mathrm{dec}}(\tau_{\mathrm{dec}})\), no
minimizer can omit a channel above the threshold while retaining one
below it. Hence every minimizer is exactly the top-\(k\) pointer subset,
up to ties above the gap.
\end{proof}

\begin{lemma}[Product-classical coarse-graining yields fragment-local omission losses]\label{lem:product_classical}
Assume that after applying the pointer coarse-graining map \(\Phi\), the
diagonal classical sector is described by \(D\) local channel records
\[
\xi=(\xi_1,\dots,\xi_D),
\qquad
\xi_i=(z_i,e_i),
\]
with full diagonal law
\[
P^{\mathrm{full}}(\xi)
=
\prod_{i=1}^D P_i(z_i,e_i).
\]
Let \(\mathcal F_I\) be an admissible pointer-adapted frame. Assume its
diagonal projected law is obtained by retaining the full local channel
law on \(I\) and replacing each omitted channel by a fixed reference
law:
\[
P^{I}(\xi)
=
\prod_{i\in I} P_i(z_i,e_i)
\prod_{i\notin I} Q_i(z_i,e_i).
\]
Assume further:
\begin{enumerate}
\item the diagonal survival identity is represented pointwise by
  \[
  \beta\,\mathcal S^{\mathrm{diag}}(\xi)
  =
  -\imath(\xi)-d(\xi)-\beta\,\delta h(\xi),
  \]
  where \(\imath\) is the pointwise mutual-information density, \(d\) is
  the pointwise reference-divergence density, and \(\delta h\) is the
  pointwise control contribution;
\item the reference divergence is taken with respect to a product
  reference law
  \[
  \tau(e)=\prod_{i=1}^D \tau_i(e_i);
  \]
\item the diagonal control contribution is additive:
  \[
  \delta h^{\mathrm{full}}(\xi)
  =
  \sum_{i=1}^D \delta h_i^{\mathrm{full}}(\xi_i),
  \]
  \[
  \delta h^{I}(\xi)
  =
  \sum_{i\in I}\delta h_i^{\mathrm{full}}(\xi_i)
  +
  \sum_{i\notin I}\delta h_i^{\mathrm{ref}}(\xi_i).
  \]
\end{enumerate}
For each channel \(i\), define the local pointwise mutual-information
density
\[
\imath_i^{\mathrm{full}}(\xi_i)
:=
\log\frac{P_i(z_i,e_i)}{P_i^Z(z_i)\,P_i^E(e_i)},
\qquad
\imath_i^{\mathrm{ref}}(\xi_i)
:=
\log\frac{Q_i(z_i,e_i)}{Q_i^Z(z_i)\,Q_i^E(e_i)},
\]
and the local reference-divergence density
\[
d_i^{\mathrm{full}}(e_i)
:=
\log\frac{P_i^E(e_i)}{\tau_i(e_i)},
\qquad
d_i^{\mathrm{ref}}(e_i)
:=
\log\frac{Q_i^E(e_i)}{\tau_i(e_i)}.
\]
Then the diagonal survival deficit is fragment-local:
\[
\beta\bigl(
\mathcal S_{\mathrm{full}}^{\mathrm{diag}}(\xi)
-
\mathcal S_{\mathcal F_I}^{\mathrm{diag}}(\xi)
\bigr)
=
\sum_{i\notin I}\lambda_i(\xi_i),
\]
where
\[
\lambda_i(\xi_i)
:=
-\bigl(\imath_i^{\mathrm{full}}(\xi_i)-\imath_i^{\mathrm{ref}}(\xi_i)\bigr)
-\bigl(d_i^{\mathrm{full}}(e_i)-d_i^{\mathrm{ref}}(e_i)\bigr)
-\beta\bigl(\delta h_i^{\mathrm{full}}(\xi_i)-\delta h_i^{\mathrm{ref}}(\xi_i)\bigr).
\]
Consequently, the centered variables
\[
L_i(\xi_i):=\lambda_i(\xi_i)-\mathbb E_\xi[\lambda_i(\xi_i)]
\]
are fragment-local omission losses in the sense needed by
Lemma~\ref{lem:fragment_orth}.
\end{lemma}

\begin{proof}
Because the full diagonal law factorizes,
\[
P^{\mathrm{full}}(\xi)
=
\prod_{i=1}^D P_i(z_i,e_i),
\quad
P_Z^{\mathrm{full}}(z)
=
\prod_{i=1}^D P_i^Z(z_i),
\quad
P_E^{\mathrm{full}}(e)
=
\prod_{i=1}^D P_i^E(e_i).
\]
Therefore the full pointwise mutual-information density is additive:
\[
\imath^{\mathrm{full}}(\xi)
:=
\log\frac{P^{\mathrm{full}}(\xi)}
{P_Z^{\mathrm{full}}(z)\,P_E^{\mathrm{full}}(e)}
=
\sum_{i=1}^D \imath_i^{\mathrm{full}}(\xi_i).
\]
The same argument for the projected law gives
\[
\imath^{I}(\xi)
=
\sum_{i\in I}\imath_i^{\mathrm{full}}(\xi_i)
+
\sum_{i\notin I}\imath_i^{\mathrm{ref}}(\xi_i).
\]

Since the reference law \(\tau\) is also a product law,
\[
d^{\mathrm{full}}(e)
:=
\log\frac{P_E^{\mathrm{full}}(e)}{\tau(e)}
=
\sum_{i=1}^D d_i^{\mathrm{full}}(e_i),
\]
and similarly
\[
d^{I}(e)
=
\sum_{i\in I} d_i^{\mathrm{full}}(e_i)
+
\sum_{i\notin I} d_i^{\mathrm{ref}}(e_i).
\]
The control contribution is additive by assumption:
\[
\delta h^{\mathrm{full}}(\xi)
=
\sum_{i=1}^D \delta h_i^{\mathrm{full}}(\xi_i),
\]
\[
\delta h^{I}(\xi)
=
\sum_{i\in I}\delta h_i^{\mathrm{full}}(\xi_i)
+
\sum_{i\notin I}\delta h_i^{\mathrm{ref}}(\xi_i).
\]

Substituting these three decompositions into the pointwise diagonal
survival identity gives
\begin{align*}
&\beta\bigl(
\mathcal S_{\mathrm{full}}^{\mathrm{diag}}(\xi)
-
\mathcal S_{\mathcal F_I}^{\mathrm{diag}}(\xi)
\bigr) \\
&=
-\bigl(\imath^{\mathrm{full}}(\xi)-\imath^{I}(\xi)\bigr)
-\bigl(d^{\mathrm{full}}(e)-d^{I}(e)\bigr)
-\beta\bigl(
\delta h^{\mathrm{full}}(\xi)-\delta h^{I}(\xi)
\bigr) \\
&=
\sum_{i\notin I}
\Bigl[
-\bigl(\imath_i^{\mathrm{full}}(\xi_i)-\imath_i^{\mathrm{ref}}(\xi_i)\bigr)
-\bigl(d_i^{\mathrm{full}}(e_i)-d_i^{\mathrm{ref}}(e_i)\bigr)
-\beta\bigl(\delta h_i^{\mathrm{full}}(\xi_i)-\delta h_i^{\mathrm{ref}}(\xi_i)\bigr)
\Bigr].
\end{align*}
This is exactly the stated fragment-local decomposition. Centering each
\(\lambda_i\) produces the centered fragment-local omission losses
\(L_i\).
\end{proof}

\begin{remark}[What Lemma~\ref{lem:product_classical} closes]
Lemma~\ref{lem:product_classical} is the first theorem-level closure of
the structure line PS2a. It does not derive the product-classical
pointer regime from the broken-phase output; it shows that once the
pointer coarse-graining reaches that regime, fragment-local omission
losses follow automatically by additivity of information densities,
reference-divergence densities, and diagonal control terms.
\end{remark}

\begin{lemma}[Orthogonality from fragment-local centered channel losses]\label{lem:fragment_orth}
Assume the environmental realization admits a pointer-channel
factorization
\[
\xi=(\xi_1,\dots,\xi_D)
\]
with independent coordinates. Suppose the diagonal omission losses of
Proposition~\ref{prop:ps2_restricted} admit a centered fragment-local
realization
\[
L_i(\xi)
=
\lambda_i(\xi_i)-\mathbb E_\xi[\lambda_i(\xi_i)],
\qquad
i=1,\dots,D,
\]
for measurable functions \(\lambda_i(\xi_i)\in L^2(\xi_i)\). Then
\[
\mathbb E_\xi[L_iL_j]=0
\qquad
\text{for } i\neq j.
\]
In particular, the pairwise orthogonality hypothesis of
Proposition~\ref{prop:ps2_restricted} is automatic under independent
pointer channels together with centered fragment-local omission losses.
\end{lemma}

\begin{proof}
Because \(L_i\) depends only on \(\xi_i\) and \(L_j\) depends only on
\(\xi_j\), the independence of \(\xi_i\) and \(\xi_j\) implies
\[
\mathbb E_\xi[L_iL_j]
=
\mathbb E_\xi[L_i]\,
\mathbb E_\xi[L_j].
\]
The centered form of \(L_i\) and \(L_j\) gives
\[
\mathbb E_\xi[L_i]=0=\mathbb E_\xi[L_j],
\]
hence \(\mathbb E_\xi[L_iL_j]=0\).
\end{proof}

\begin{remark}[Why the orthogonality hypothesis is natural in the pointer regime]
Lemma~\ref{lem:fragment_orth} explains why the first hypothesis of
Proposition~\ref{prop:ps2_restricted} is not an arbitrary add-on. On
the T-DOME side, the computational-ceiling analysis already treats the
active channels as independent rate carriers. On the HAFF/Quantum
Darwinism side, accessible classical records are carried by independent
environmental fragments. Once the diagonal omission loss localizes to
fragment-wise centered channel contributions after coarse-graining,
pairwise orthogonality becomes a probability identity rather than an
extra dynamical postulate.

What remains open is exactly the next step: deriving that fragment-local
centered realization, together with the uniform coherence bound, from
the full T-DOME survival identity in a model-independent way.
\end{remark}

\begin{lemma}[Semigroup-decay coherence estimate]\label{lem:semigroup_coherence}
Let
\[
\mathbb E_{\mathrm{ptr}}:\mathcal A_U\to \mathcal A_{\mathrm{ptr}}
\]
be a normal conditional expectation onto the pointer algebra, and write
\[
Q_{\mathrm{off}}:=\mathrm{id}-\mathbb E_{\mathrm{ptr}}.
\]
For each admissible pointer-adapted frame \(\mathcal F_I\) and
environmental realization \(\xi\), let \(X_\xi\) denote the relevant
state or record datum entering the survival identity, and decompose the
coherence remainder as
\[
C_I(\xi)
=
C_I^{(I)}(\xi)+C_I^{(D)}(\xi)+C_I^{(H)}(\xi),
\]
where the three terms correspond respectively to the off-diagonal
contributions of
\[
-\Delta I(S{:}E),
\qquad
-\Delta D_{\mathrm{KL}},
\qquad
-\beta\,\Delta\langle H_{\mathrm{ctrl}}\rangle.
\]
Assume:
\begin{enumerate}
\item \textbf{Off-sector semigroup decay:} there exists a non-increasing
  function \(g:[0,\infty)\to [0,\infty)\) with \(g(\tau)\to 0\) such
  that
  \[
  \sup_{t\ge \tau}
  \|Q_{\mathrm{off}}\mathcal E_t(X)\|
  \le
  g(\tau)\,\|Q_{\mathrm{off}}X\|
  \]
  for every relevant datum \(X\);
\item \textbf{Uniform off-sector size:}
  \[
  \sup_\xi \|Q_{\mathrm{off}}X_\xi\|
  \le
  B_{\mathrm{off}}
  <
  \infty;
  \]
\item \textbf{Termwise continuity bounds:} there exist constants
  \(L_I,L_D,L_H\ge 0\) such that for every admissible \(I\) and every
  \(\xi\),
  \[
  |C_I^{(I)}(\xi)|
  \le
  L_I\sup_{t\ge\tau_{\mathrm{dec}}}
  \|Q_{\mathrm{off}}\mathcal E_t(X_\xi)\|,
  \]
  \[
  |C_I^{(D)}(\xi)|
  \le
  L_D\sup_{t\ge\tau_{\mathrm{dec}}}
  \|Q_{\mathrm{off}}\mathcal E_t(X_\xi)\|,
  \]
  \[
  |C_I^{(H)}(\xi)|
  \le
  L_H\sup_{t\ge\tau_{\mathrm{dec}}}
  \|Q_{\mathrm{off}}\mathcal E_t(X_\xi)\|.
  \]
\end{enumerate}
Then for every admissible \(I\),
\[
\|C_I\|_{L^\infty(\xi)}
\le
\eta_{\mathrm{dec}}(\tau_{\mathrm{dec}})
:=
(L_I+L_D+L_H)\,B_{\mathrm{off}}\,g(\tau_{\mathrm{dec}}).
\]
In particular, the uniform coherence bound required in
Proposition~\ref{prop:ps2_restricted} holds with this choice of
\(\eta_{\mathrm{dec}}\).
\end{lemma}

\begin{proof}
By the triangle inequality,
\[
|C_I(\xi)|
\le
|C_I^{(I)}(\xi)|+|C_I^{(D)}(\xi)|+|C_I^{(H)}(\xi)|.
\]
Applying the termwise continuity bounds gives
\[
|C_I(\xi)|
\le
(L_I+L_D+L_H)\,
\sup_{t\ge\tau_{\mathrm{dec}}}
\|Q_{\mathrm{off}}\mathcal E_t(X_\xi)\|.
\]
The off-sector semigroup decay assumption then yields
\[
|C_I(\xi)|
\le
(L_I+L_D+L_H)\,
g(\tau_{\mathrm{dec}})\,
\|Q_{\mathrm{off}}X_\xi\|.
\]
Using the uniform off-sector bound,
\[
|C_I(\xi)|
\le
(L_I+L_D+L_H)\,
B_{\mathrm{off}}\,
g(\tau_{\mathrm{dec}}).
\]
Taking the essential supremum over \(\xi\) gives the stated
\(L^\infty\) bound.
\end{proof}

\begin{remark}[What Lemma~\ref{lem:semigroup_coherence} closes]
Lemma~\ref{lem:semigroup_coherence} is the abstract translation step
from standard decoherence-semigroup suppression of the off-pointer
sector to the exact \(\eta_{\mathrm{dec}}\) object used in
Proposition~\ref{prop:ps2_restricted}. It does not yet identify the
best norm or derive the constants \(L_I,L_D,L_H\) from the full
T-DOME survival identity; those remain the true analytic content.

In concrete pointer-diagonal control models, one often has
\(L_H=0\), so only the mutual-information and relative-entropy terms
contribute to the coherence remainder.
\end{remark}

\begin{corollary}[Hilbert--Schmidt continuity package on a uniformly faithful pointer support]\label{cor:hs_package}
Assume that, after pointer coarse-graining, the relevant joint
system--environment states are supported on a finite-dimensional space
\[
\mathcal H_S\otimes \mathcal H_E
\]
with
\[
d_S:=\dim\mathcal H_S,
\qquad
d_E:=\dim\mathcal H_E.
\]
Assume there exists \(m_*>0\) such that every relevant joint state
\(\rho_{SE}\), every relevant marginal state \(\rho_S,\rho_E\), and the
fixed reference state \(\tau_E\) entering the T-DOME
\(\Delta D_{\mathrm{KL}}\) term satisfy
\[
\lambda_{\min}(\rho_{SE}),
\lambda_{\min}(\rho_S),
\lambda_{\min}(\rho_E),
\lambda_{\min}(\tau_E)
\ge
m_*.
\]
Write \(\|\cdot\|_2\) for the Hilbert--Schmidt norm.

Then:
\begin{enumerate}
\item for any two relevant joint states \(\rho_{SE},\sigma_{SE}\),
  \[
  \bigl|I(\rho_{SE})-I(\sigma_{SE})\bigr|
  \le
  3\sqrt{d_Sd_E}\,\bigl(|\log m_*|+1\bigr)\,
  \|\rho_{SE}-\sigma_{SE}\|_2;
  \]
\item for any two relevant environment states \(\rho_E,\sigma_E\),
  \[
  \bigl|D_{\mathrm{KL}}(\rho_E\|\tau_E)
  -D_{\mathrm{KL}}(\sigma_E\|\tau_E)\bigr|
  \le
  \sqrt{d_E}\,\bigl(2|\log m_*|+1\bigr)\,
  \|\rho_E-\sigma_E\|_2;
  \]
  \item for any self-adjoint control Hamiltonian
  \(H_{\mathrm{ctrl}}\in\mathcal B(\mathcal H_S)\),
  \[
  \bigl|\operatorname{Tr}[H_{\mathrm{ctrl}}(\rho_S-\sigma_S)]\bigr|
  \le
  \|H_{\mathrm{ctrl}}\|_2\,
  \|\rho_S-\sigma_S\|_2.
  \]
\end{enumerate}
Consequently, in this restricted finite-dimensional faithful regime, one
may take Hilbert--Schmidt continuity constants
\[
L_I^{\mathrm{HS}}
:=
3\sqrt{d_Sd_E}\,\bigl(|\log m_*|+1\bigr),
\]
\[
L_D^{\mathrm{HS}}
:=
\sqrt{d_E}\,\bigl(2|\log m_*|+1\bigr),
\]
\[
L_H^{\mathrm{HS}}
:=
\|H_{\mathrm{ctrl}}\|_2
\]
in Lemma~\ref{lem:semigroup_coherence}, after translating the relevant
off-sector perturbations into the corresponding Hilbert--Schmidt norms.
\end{corollary}

\begin{proof}
Set
\[
S(\rho):=-\operatorname{Tr}(\rho\log\rho).
\]
On the convex set of density operators with spectrum contained in
\([m_*,1]\), the Fr\'echet derivative is
\[
DS_\rho(H)=-\operatorname{Tr}\bigl(H(\log\rho+\mathbf 1)\bigr).
\]
Hence
\[
|DS_\rho(H)|
\le
\|H\|_2\,\|\log\rho+\mathbf 1\|_2
\le
\sqrt{d}\,\bigl(|\log m_*|+1\bigr)\|H\|_2
\]
for \(d=\dim\operatorname{supp}\rho\). Applying the mean-value theorem
along the line segment between two states gives
\[
|S(\rho)-S(\sigma)|
\le
\sqrt{d}\,\bigl(|\log m_*|+1\bigr)\|\rho-\sigma\|_2.
\]

For the mutual-information term, use
\[
I(\rho_{SE})=S(\rho_S)+S(\rho_E)-S(\rho_{SE}).
\]
The partial trace satisfies the Hilbert--Schmidt bound
\[
\|\operatorname{Tr}_E A\|_2\le \sqrt{d_E}\,\|A\|_2,
\qquad
\|\operatorname{Tr}_S A\|_2\le \sqrt{d_S}\,\|A\|_2.
\]
Therefore
\[
|S(\rho_S)-S(\sigma_S)|
\le
\sqrt{d_S}\,\bigl(|\log m_*|+1\bigr)\|\rho_S-\sigma_S\|_2
\le
\sqrt{d_Sd_E}\,\bigl(|\log m_*|+1\bigr)\|\rho_{SE}-\sigma_{SE}\|_2,
\]
and similarly for the environment marginal. For the joint entropy term,
\[
|S(\rho_{SE})-S(\sigma_{SE})|
\le
\sqrt{d_Sd_E}\,\bigl(|\log m_*|+1\bigr)\|\rho_{SE}-\sigma_{SE}\|_2.
\]
Summing the three contributions proves part~(1).

For the relative-entropy term with fixed faithful reference \(\tau_E\),
write
\[
D_{\mathrm{KL}}(\rho_E\|\tau_E)
=
\operatorname{Tr}(\rho_E\log\rho_E)-\operatorname{Tr}(\rho_E\log\tau_E).
\]
Hence
\begin{align*}
&\bigl|D_{\mathrm{KL}}(\rho_E\|\tau_E)
-D_{\mathrm{KL}}(\sigma_E\|\tau_E)\bigr| \\
&\le
\bigl|S(\rho_E)-S(\sigma_E)\bigr|
+
\bigl|\operatorname{Tr}\bigl((\rho_E-\sigma_E)\log\tau_E\bigr)\bigr|.
\end{align*}
The entropy term is bounded by
\[
\sqrt{d_E}\,\bigl(|\log m_*|+1\bigr)\|\rho_E-\sigma_E\|_2,
\]
while
\[
\|\log\tau_E\|_2\le \sqrt{d_E}\,|\log m_*|.
\]
Cauchy--Schwarz therefore gives part~(2).

Part~(3) is immediate from the Hilbert--Schmidt Cauchy--Schwarz
inequality.
\end{proof}

\begin{remark}[Status of Corollary~\ref{cor:hs_package}]
Corollary~\ref{cor:hs_package} does not close the full estimate track in
general Type~III form. It closes the first finite-dimensional faithful
pointer-support regime, where the Hilbert--Schmidt norm is natural and
all three continuity constants become explicit. For the restricted
Bridge~1 target developed in the present note, that is enough: the
estimate line is closed at theorem level. The further question of
whether an equally clean package survives beyond that regime, or must be
reformulated in a predual norm, belongs to a later strengthening rather
than to the current closure claim.
\end{remark}

\begin{corollary}[Restricted PS2 closure in the product-classical faithful pointer regime]\label{cor:ps2_closure}
Assume:
\begin{enumerate}
\item the hypotheses of Lemma~\ref{lem:product_classical};
\item the independent pointer-channel factorization hypothesis of
  Lemma~\ref{lem:fragment_orth};
\item the off-sector semigroup decay and uniform off-sector size
  hypotheses of Lemma~\ref{lem:semigroup_coherence};
\item the finite-dimensional uniformly faithful pointer-support
  hypotheses of Corollary~\ref{cor:hs_package};
\item the centered fragment-local omission losses from
  Lemma~\ref{lem:product_classical} satisfy
  \[
  \Lambda_k
  :=
  \sup_{|I|=k}
  \sum_{i\notin I}\|L_i\|_{L^2(\xi)}
  <
  \infty.
  \]
\end{enumerate}
Define
\[
w_i:=\|L_i\|_{L^2(\xi)}^2,
\]
\[
\eta_{\mathrm{dec}}^{\mathrm{HS}}(\tau_{\mathrm{dec}})
:=
\bigl(
L_I^{\mathrm{HS}}+L_D^{\mathrm{HS}}+L_H^{\mathrm{HS}}
\bigr)
B_{\mathrm{off}}\,g(\tau_{\mathrm{dec}}),
\]
and
\[
\varepsilon_{\mathrm{dec}}^{\mathrm{HS}}(\tau_{\mathrm{dec}})
:=
2\Lambda_k\,\eta_{\mathrm{dec}}^{\mathrm{HS}}(\tau_{\mathrm{dec}})
+
\bigl(\eta_{\mathrm{dec}}^{\mathrm{HS}}(\tau_{\mathrm{dec}})\bigr)^2.
\]
Then every admissible pointer-adapted frame \(\mathcal F_I\) satisfies
\[
D(\mathcal F_I)
=
\sum_{i\notin I} w_i + R_I,
\qquad
|R_I|
\le
\varepsilon_{\mathrm{dec}}^{\mathrm{HS}}(\tau_{\mathrm{dec}}).
\]
Consequently, if
\[
\Delta_k:=w_{(k)}-w_{(k+1)}
>
2\varepsilon_{\mathrm{dec}}^{\mathrm{HS}}(\tau_{\mathrm{dec}}),
\]
then every minimizer of \(D(\mathcal F_I)\) retains exactly the top-\(k\)
pointer channels, up to ties above the gap.
\end{corollary}

\begin{proof}
Lemma~\ref{lem:product_classical} produces centered fragment-local
omission losses \(L_i\). Lemma~\ref{lem:fragment_orth} upgrades the
independent pointer-channel factorization to the pairwise orthogonality
required in Proposition~\ref{prop:ps2_restricted}. By
Corollary~\ref{cor:hs_package}, the continuity constants appearing in
Lemma~\ref{lem:semigroup_coherence} may be chosen as
\(L_I^{\mathrm{HS}},L_D^{\mathrm{HS}},L_H^{\mathrm{HS}}\), so that
Lemma~\ref{lem:semigroup_coherence} yields the uniform coherence bound
\[
\|C_I\|_{L^\infty(\xi)}
\le
\eta_{\mathrm{dec}}^{\mathrm{HS}}(\tau_{\mathrm{dec}}).
\]
The finiteness hypothesis gives the remaining size bound
\(\Lambda_k<\infty\). Proposition~\ref{prop:ps2_restricted} therefore
applies and gives exactly the stated approximate separability law and
top-\(k\) selection criterion.
\end{proof}

\begin{corollary}[Pure-dephasing envelope]\label{cor:ps2_pure_dephasing}
Assume the hypotheses of Proposition~\ref{prop:ps2_restricted}, and
suppose in addition that the coherence remainder is controlled by the
omitted pointer-channel coherences in the form
\[
|C_I(\xi,t)|
\le
\sum_{i\notin I} a_i\,|p_i(t)|
\qquad
\text{for all } t\ge \tau_{\mathrm{dec}},
\]
where \(a_i\ge 0\) are constants and each \(p_i(t)\) is the exact
Lorentz--Drude pure-dephasing decoherence function from
T-DOME~I~\cite{TDOMEI,BreuerPetruccione2002},
\[
\gamma_i<4\lambda_i,
\qquad
\Omega_i
:=
\tfrac{1}{2}\sqrt{4\lambda_i\gamma_i-\gamma_i^2}
\in \mathbb R_{>0},
\]
and
\[
p_i(t)
=
e^{-\gamma_i t/2}\left[
\cos(\Omega_i t)
+
\frac{\gamma_i}{2\Omega_i}\sin(\Omega_i t)
\right].
\]
Define
\[
R_i
:=
\sqrt{1+\left(\frac{\gamma_i}{2\Omega_i}\right)^2},
\qquad
A_k
:=
\max_{|I|=k}\sum_{i\notin I} a_iR_i.
\]
Then one may take
\[
\eta_{\mathrm{dec}}(\tau_{\mathrm{dec}})
=
A_k\,e^{-\gamma_* \tau_{\mathrm{dec}}/2},
\qquad
\gamma_*
:=
\min_{1\le i\le D}\gamma_i,
\]
and hence
\[
\varepsilon_{\mathrm{dec}}(\tau_{\mathrm{dec}})
\le
2\Lambda_k A_k\,e^{-\gamma_* \tau_{\mathrm{dec}}/2}
+
A_k^2\,e^{-\gamma_* \tau_{\mathrm{dec}}}.
\]
In particular, the distortion remainder decays exponentially with the
decoherence depth.
\end{corollary}

\begin{proof}
For each channel,
\[
|p_i(t)|
\le
e^{-\gamma_i t/2}
\sqrt{1+\left(\frac{\gamma_i}{2\Omega_i}\right)^2}
=
R_i\,e^{-\gamma_i t/2}
\le
R_i\,e^{-\gamma_* t/2}
\]
because \(|a\cos x+b\sin x|\le \sqrt{a^2+b^2}\). Therefore
\[
|C_I(\xi,t)|
\le
\sum_{i\notin I} a_iR_i\,e^{-\gamma_* t/2}
\le
A_k\,e^{-\gamma_* \tau_{\mathrm{dec}}/2}
\]
for all \(t\ge \tau_{\mathrm{dec}}\). This identifies a valid choice of
\(\eta_{\mathrm{dec}}(\tau_{\mathrm{dec}})\). The stated bound on
\(\varepsilon_{\mathrm{dec}}(\tau_{\mathrm{dec}})\) is then immediate
from Proposition~\ref{prop:ps2_restricted}.
\end{proof}

\paragraph{A bookkeeping split for the model-independent PS2 programme.}
The restricted results above close the PS2 package only after explicit
regime assumptions have been imposed. For the later model-independent
programme, one further lemma is useful already now: it fixes the exact
place where the off-sector estimates enter the survival deficit.

For \(X\in\{\mathrm{full},\mathcal F_I\}\), write the survival identity
as
\[
\beta\,\mathcal S_X
=
-\Delta I_X
-\Delta D_X
-\beta\,\Delta H_X,
\]
where
\[
\Delta D_X:=\Delta D_{\mathrm{KL},X},
\qquad
\Delta H_X:=\Delta\langle H_{\mathrm{ctrl}}\rangle_X.
\]
Assume each of the three terms admits a pointer-diagonal/off-sector
decomposition,
\[
\Delta I_X
:=
\Delta I_X^{\mathrm{diag}}
+
\Delta I_X^{\mathrm{off}},
\]
\[
\Delta D_X
:=
\Delta D_X^{\mathrm{diag}}
+
\Delta D_X^{\mathrm{off}},
\]
\[
\Delta H_X
:=
\Delta H_X^{\mathrm{diag}}
+
\Delta H_X^{\mathrm{off}}.
\]
Define the pointer-diagonal survival functional and off-sector remainder
by
\[
\beta\,\mathcal S_X^{\mathrm{diag}}
:=
-\Delta I_X^{\mathrm{diag}}
-\Delta D_X^{\mathrm{diag}}
-\beta\,\Delta H_X^{\mathrm{diag}},
\]
\[
R_X^{\mathrm{off}}
:=
-\Delta I_X^{\mathrm{off}}
-\Delta D_X^{\mathrm{off}}
-\beta\,\Delta H_X^{\mathrm{off}}.
\]

\begin{lemma}[Diagonal/off-sector split of the survival deficit]\label{lem:diag_off_split}
Let \(\mathbb E_{\mathrm{ptr}}:\mathcal A_U\to \mathcal A_{\mathrm{ptr}}\)
be the conditional expectation onto the pointer algebra selected by the
broken phase, and assume the above termwise decompositions are computed
with respect to \(\mathbb E_{\mathrm{ptr}}\). Then for every admissible
pointer-adapted frame \(\mathcal F_I\),
\[
\beta\bigl(
\mathcal S_{\mathrm{full}}
-
\mathcal S_{\mathcal F_I}
\bigr)
=
\beta\bigl(
\mathcal S_{\mathrm{full}}^{\mathrm{diag}}
-
\mathcal S_{\mathcal F_I}^{\mathrm{diag}}
\bigr)
+
R_I^{\mathrm{off}},
\]
where
\[
R_I^{\mathrm{off}}
:=
R_{\mathrm{full}}^{\mathrm{off}}
-
R_{\mathcal F_I}^{\mathrm{off}}.
\]
Equivalently,
\begin{align*}
R_I^{\mathrm{off}}
=&
-\bigl(
\Delta I_{\mathrm{full}}^{\mathrm{off}}
-
\Delta I_{\mathcal F_I}^{\mathrm{off}}
\bigr) \\
&
-\bigl(
\Delta D_{\mathrm{full}}^{\mathrm{off}}
-
\Delta D_{\mathcal F_I}^{\mathrm{off}}
\bigr) \\
&
-\beta\bigl(
\Delta H_{\mathrm{full}}^{\mathrm{off}}
-
\Delta H_{\mathcal F_I}^{\mathrm{off}}
\bigr).
\end{align*}
\end{lemma}

\begin{proof}
Write the survival identity separately for the full process and the
frame-adapted process \(\mathcal F_I\), insert the termwise
diagonal/off-sector decompositions into both identities, and subtract.
The pointer-diagonal contributions group into
\[
\beta\bigl(
\mathcal S_{\mathrm{full}}^{\mathrm{diag}}
-
\mathcal S_{\mathcal F_I}^{\mathrm{diag}}
\bigr),
\]
while the off-sector contributions group into
\[
R_{\mathrm{full}}^{\mathrm{off}}
-
R_{\mathcal F_I}^{\mathrm{off}}.
\]
This is exactly the stated identity.
\end{proof}

\begin{remark}[What this does and does not close]
Section~\ref{sec:ps2model} is a restricted theorem-level closure of PS2,
not a general one. What is now closed is the implication
\[
\begin{aligned}
\text{uniformly small coherence remainder}
&\Longrightarrow
\text{approximate separability}
\\
&\Longrightarrow
\text{top-}k \text{ pointer selection}.
\end{aligned}
\]
Corollary~\ref{cor:ps2_closure} packages this as a single restricted
selection theorem in the product-classical and finite-dimensional
faithful pointer regime, while
Corollary~\ref{cor:ps2_pure_dephasing} gives a model-specific
exponential envelope. Lemma~\ref{lem:diag_off_split} does not enlarge
that closure claim; it only fixes the bookkeeping split showing that
the already-controlled semigroup-decay terms belong to the off-sector
remainder. What remains open is the model-independent derivation of the
regime assumptions, and in particular control of the diagonal product
defect, from the full T-DOME survival identity and the broken-phase
output.
\end{remark}

\paragraph{A conditional model-independent foreground-background package.}
The bookkeeping split also makes it possible to state an honest
conditional version of the next model-independent target. For
budget-feasible \(I\) with \(|I|=k^*\), define the diagonal-tail
survival deficit by
\[
\Delta_I^{\mathrm{diag}}(\xi,t)
:=
\beta\bigl(
\mathcal S_{\mathrm{full}}^{\mathrm{diag}}(\xi,t)
-
\mathcal S_{\mathcal F_I}^{\mathrm{diag}}(\xi,t)
\bigr),
\]
the corresponding tail distortion by
\[
D_\tau^{\mathrm{diag}}(I)
:=
\sup_{t\ge\tau}
\mathbb E_{P_t^{\mathrm{diag}}}\!\Bigl[
\bigl(\Delta_I^{\mathrm{diag}}(\xi,t)\bigr)^2
\Bigr],
\]
and for an exchange pair
\[
\mathfrak P:=I^*\times (I^*)^{\mathrm c},
\qquad
N_{\mathrm{ex}}:=|\mathfrak P|=k^*(D-k^*),
\]
define
\[
G_{a\leftrightarrow b}^{\mathrm{diag}}(\xi,t)
:=
\beta\bigl(
\mathcal S_{\mathcal F_{I_{a\leftrightarrow b}}}^{\mathrm{diag}}(\xi,t)
-
\mathcal S_{\mathcal F_{I^*}}^{\mathrm{diag}}(\xi,t)
\bigr),
\]
\[
\gamma_{\mathrm{fb}}^{\mathrm{diag}}(\xi,t)
:=
\beta\bigl(
\mathcal S_{\mathrm{full}}^{\mathrm{diag}}(\xi,t)
-
\mathcal S_{\mathrm{ff}}^{\mathrm{diag}}(u,t)
-
\mathcal S_{\mathrm{bb}}^{\mathrm{diag}}(v,t)
\bigr),
\]
\[
\overline{\Gamma}_{\mathrm{fb}}^{\mathrm{diag}}(t)
:=
\mathbb E_{P_t^{\mathrm{diag}}}\!\bigl[
\gamma_{\mathrm{fb}}^{\mathrm{diag}}(\xi,t)
\bigr].
\]

\begin{proposition}[Conditional diagonal foreground-background coupling suppression]\label{prop:conditional_b2a}
Let \(I^*\) be the broken-phase foreground set. Assume:
\begin{enumerate}
\item \textbf{Approximate tail optimality.} There exists
  \(r_{\mathrm{opt}}(\tau)\to 0\) such that
  \[
  D_\tau^{\mathrm{diag}}(I^*)
  \le
  D_\tau^{\mathrm{diag}}(I)+r_{\mathrm{opt}}(\tau)
  \qquad
  \text{for all } |I|=k^*.
  \]
\item \textbf{Exchange cover and tail persistence.} There exist
  \(c_{\mathrm{cov}}>0\), \(M_\gamma<\infty\),
  \(e_{\mathrm{cov}}(\tau)\to 0\), and
  \(r_{\mathrm{pers}}(\tau)\to 0\) such that:
  \[
  G_{a\leftrightarrow b}^{\mathrm{diag}}(\xi,t)\,
  \gamma_{\mathrm{fb}}^{\mathrm{diag}}(\xi,t)\ge 0
  \qquad
  \text{for all } (a,b)\in\mathfrak P,\ t\ge\tau,
  \]
  \[
  \frac{1}{N_{\mathrm{ex}}}
  \sum_{(a,b)\in\mathfrak P}
  \bigl|
  G_{a\leftrightarrow b}^{\mathrm{diag}}(\xi,t)
  \bigr|
  \ge
  c_{\mathrm{cov}}
  \bigl|
  \gamma_{\mathrm{fb}}^{\mathrm{diag}}(\xi,t)
  \bigr|
  -
  e_{\mathrm{cov}}(\tau)
  \quad
  \text{for a.e. } \xi,\ t\ge\tau,
  \]
  \[
  \sup_{t\ge\tau}
  \mathbb E\!\bigl[
  |\gamma_{\mathrm{fb}}^{\mathrm{diag}}(\cdot,t)|
  \bigr]
  \le
  M_\gamma,
  \]
  and for
  \[
  \Gamma_{a\leftrightarrow b}^{\mathrm{ex}}(t)
  :=
  \mathbb E\!\bigl[
  G_{a\leftrightarrow b}^{\mathrm{diag}}(\cdot,t)\,
  \gamma_{\mathrm{fb}}^{\mathrm{diag}}(\cdot,t)
  \bigr]
  \]
  one has
  \[
  \sup_{(a,b)\in\mathfrak P}
  \sup_{s,t\ge\tau}
  \bigl|
  \Gamma_{a\leftrightarrow b}^{\mathrm{ex}}(t)
  -
  \Gamma_{a\leftrightarrow b}^{\mathrm{ex}}(s)
  \bigr|
  \le
  r_{\mathrm{pers}}(\tau).
  \]
\item \textbf{Tail-passive background survival.} There exist
  \(M_{\mathrm{ex}}(\tau)<\infty\) and
  \(\eta_{\mathrm{bb}}(\tau)\to 0\) such that
  \[
  \sup_{(a,b)\in\mathfrak P}\sup_{t\ge\tau}
  \bigl\|
  G_{a\leftrightarrow b}^{\mathrm{diag}}(\cdot,t)
  \bigr\|_{L^1}
  \le
  M_{\mathrm{ex}}(\tau),
  \]
  \[
  \sup_{t\ge\tau}
  \bigl\|
  \mathcal S_{\mathrm{bb}}^{\mathrm{diag}}(\cdot,t)
  \bigr\|_{L^\infty}
  \le
  \eta_{\mathrm{bb}}(\tau).
  \]
\item \textbf{Exchange efficiency gap.} There exist
  \(M_G(\tau)<\infty\), \(m_\gamma(\tau)>0\), and a constant
  \(\delta_Q>0\) such that
  \[
  \sup_{(a,b)\in\mathfrak P}\sup_{t\ge\tau}
  \bigl\|
  G_{a\leftrightarrow b}^{\mathrm{diag}}(\cdot,t)
  \bigr\|_{L^\infty}
  \le
  M_G(\tau),
  \]
  \[
  \bigl|
  \gamma_{\mathrm{fb}}^{\mathrm{diag}}(\xi,t)
  \bigr|
  \ge
  m_\gamma(\tau)
  \quad
  \text{whenever }
  G_{a\leftrightarrow b}^{\mathrm{diag}}(\xi,t)\neq 0,
  \]
  and the efficiency ratio
  \[
  \kappa_Q(\tau):=\frac{M_G(\tau)}{m_\gamma(\tau)}
  \]
  satisfies
  \[
  \kappa_Q(\tau)\le 2-\delta_Q
  \qquad
  \text{for all sufficiently large } \tau.
  \]
\end{enumerate}
Then there exists a function \(\varepsilon_\Gamma(\tau)\to 0\) such
that
\[
\sup_{t\ge\tau}
\bigl|
\overline{\Gamma}_{\mathrm{fb}}^{\mathrm{diag}}(t)
\bigr|
\le
\varepsilon_\Gamma(\tau)
\]
for all sufficiently large \(\tau\). One may take
\[
\varepsilon_\Gamma(\tau)
:=
\left[
\frac{N_{\mathrm{ex}}}{\delta_Q c_{\mathrm{cov}}}
\left(
r_{\mathrm{opt}}(\tau)
+
\frac{2e_{\mathrm{cov}}(\tau)}{N_{\mathrm{ex}}}M_\gamma
+
2r_{\mathrm{pers}}(\tau)
+
2\beta M_{\mathrm{ex}}(\tau)\eta_{\mathrm{bb}}(\tau)
\right)
\right]^{1/2}.
\]
\end{proposition}

\begin{proof}
The exact exchange expansion gives, for every exchange pair,
\[
D_\tau^{\mathrm{diag}}(I_{a\leftrightarrow b})
-
D_\tau^{\mathrm{diag}}(I^*)
\le
Q_{a\leftrightarrow b}^\tau
+
2B_{a\leftrightarrow b}^\tau
-
2A_{a\leftrightarrow b}^\tau,
\]
where
\[
Q_{a\leftrightarrow b}^\tau
:=
\sup_{t\ge\tau}
\mathbb E\!\Bigl[
\bigl(G_{a\leftrightarrow b}^{\mathrm{diag}}(\cdot,t)\bigr)^2
\Bigr],
\]
\[
B_{a\leftrightarrow b}^\tau
:=
\sup_{t\ge\tau}
\left|
\beta\,
\mathbb E\!\Bigl[
G_{a\leftrightarrow b}^{\mathrm{diag}}(\cdot,t)\,
\mathcal S_{\mathrm{bb}}^{\mathrm{diag}}(\cdot,t)
\Bigr]
\right|,
\]
\[
A_{a\leftrightarrow b}^\tau
:=
\inf_{t\ge\tau}
\Gamma_{a\leftrightarrow b}^{\mathrm{ex}}(t).
\]
If
\[
\sup_{t\ge\tau}
\bigl|
\overline{\Gamma}_{\mathrm{fb}}^{\mathrm{diag}}(t)
\bigr|
>
\varepsilon,
\]
then the exchange-cover and tail-persistence assumptions imply the
existence of \((a,b)\in\mathfrak P\) with
\[
A_{a\leftrightarrow b}^\tau
\ge
\frac{c_{\mathrm{cov}}}{N_{\mathrm{ex}}}\,\varepsilon^2
-
\frac{e_{\mathrm{cov}}(\tau)}{N_{\mathrm{ex}}}M_\gamma
-
r_{\mathrm{pers}}(\tau).
\]
The tail-passive background assumption gives
\[
2B_{a\leftrightarrow b}^\tau
\le
2\beta M_{\mathrm{ex}}(\tau)\eta_{\mathrm{bb}}(\tau).
\]
The exchange-efficiency gap implies
\[
Q_{a\leftrightarrow b}^\tau
\le
(2-\delta_Q)\,
\sup_{t\ge\tau}\Gamma_{a\leftrightarrow b}^{\mathrm{ex}}(t)
\le
(2-\delta_Q)\bigl(
A_{a\leftrightarrow b}^\tau+r_{\mathrm{pers}}(\tau)
\bigr).
\]
Substituting these estimates into the exchange bound gives
\[
D_\tau^{\mathrm{diag}}(I_{a\leftrightarrow b})
-
D_\tau^{\mathrm{diag}}(I^*)
\le
-
\frac{\delta_Q c_{\mathrm{cov}}}{N_{\mathrm{ex}}}\,\varepsilon^2
+
\frac{2e_{\mathrm{cov}}(\tau)}{N_{\mathrm{ex}}}M_\gamma
+
2r_{\mathrm{pers}}(\tau)
+
2\beta M_{\mathrm{ex}}(\tau)\eta_{\mathrm{bb}}(\tau).
\]
Approximate tail optimality of \(I^*\) yields
\[
-r_{\mathrm{opt}}(\tau)
\le
D_\tau^{\mathrm{diag}}(I_{a\leftrightarrow b})
-
D_\tau^{\mathrm{diag}}(I^*).
\]
Combining the two inequalities gives the stated bound on
\(\varepsilon\), hence on
\(\sup_{t\ge\tau}|\overline{\Gamma}_{\mathrm{fb}}^{\mathrm{diag}}(t)|\).
\end{proof}

\begin{corollary}[Conditional foreground-background mutual-information decay]\label{cor:conditional_b2b}
Assume the hypotheses of Proposition~\ref{prop:conditional_b2a}, and in
addition suppose that the pointwise diagonal foreground-background defect
admits the block-local decomposition
\[
\gamma_{\mathrm{fb}}^{\mathrm{diag}}(\xi,t)
=
-\imath_{\mathrm{fb}}(\xi,t)
+
e_{\mathrm{fb}}^D(\xi,t)
+
e_{\mathrm{fb}}^H(\xi,t),
\]
where
\[
\imath_{\mathrm{fb}}(\xi,t)
:=
\log\frac{P_t^{\mathrm{diag}}(u,v)}
{P_t^{\mathrm{diag}}(u)P_t^{\mathrm{diag}}(v)}
\]
and the averaged locality defects satisfy
\[
\sup_{t\ge\tau}
\bigl|
\overline{E}_{\mathrm{fb}}^D(t)
\bigr|
\le
\varepsilon_D(\tau),
\qquad
\sup_{t\ge\tau}
\bigl|
\overline{E}_{\mathrm{fb}}^H(t)
\bigr|
\le
\varepsilon_H(\tau),
\]
with \(\varepsilon_D(\tau),\varepsilon_H(\tau)\to 0\). Then
\[
\sup_{t\ge\tau}
I_t(u:v)
\le
\varepsilon_\Gamma(\tau)
+
\varepsilon_D(\tau)
+
\varepsilon_H(\tau),
\]
hence \(I_t(u:v)\to 0\) uniformly on the tail.
\end{corollary}

\begin{proof}
Averaging the assumed pointwise decomposition gives
\[
\overline{\Gamma}_{\mathrm{fb}}^{\mathrm{diag}}(t)
=
-I_t(u:v)
+
\overline{E}_{\mathrm{fb}}^D(t)
+
\overline{E}_{\mathrm{fb}}^H(t).
\]
Taking absolute values and applying
Proposition~\ref{prop:conditional_b2a} gives the stated estimate.
\end{proof}

\begin{corollary}[Conditional full diagonal product-defect closure]\label{cor:conditional_full_prod}
Fix orderings
\[
B_{\mathrm{fg}}=(i_1,\dots,i_{k^*}),
\qquad
B_{\mathrm{bg}}=(j_1,\dots,j_{D-k^*})
\]
of the foreground and background indices. For an ordered block
\(B=(\ell_1,\dots,\ell_r)\), write
\[
X_{B,\le m}:=(\xi_{\ell_1},\dots,\xi_{\ell_m}),
\qquad
I_{B,m}(t):=I_t\!\bigl(X_{B,\le m}:\xi_{\ell_{m+1}}\bigr).
\]
Assume the hypotheses of Corollary~\ref{cor:conditional_b2b}, and in
addition suppose that for each ordered block
\(B\in\{B_{\mathrm{fg}},B_{\mathrm{bg}}\}\) and each
\(m=1,\dots,|B|-1\) there is a pointwise diagonal prefix defect
\[
\gamma_{B,m}^{\mathrm{diag}}(x_{B,\le m+1},t)
=
-\imath_{B,m}(x_{B,\le m+1},t)
+
e_{B,m}^D(x_{B,\le m+1},t)
+
e_{B,m}^H(x_{B,\le m+1},t),
\]
where
\[
\imath_{B,m}(x_{B,\le m+1},t)
:=
\log\frac{
P_{B,t}^{\le m+1}(x_{B,\le m+1})
}{
P_{B,t}^{\le m}(x_{B,\le m})\,P_{\ell_{m+1},t}(x_{\ell_{m+1}})
},
\]
and the averaged defects satisfy
\[
\sup_{t\ge\tau}
\bigl|
\overline{\Gamma}_{B,m}^{\mathrm{diag}}(t)
\bigr|
\le
\varepsilon_{B,m}^{\Gamma}(\tau),
\]
\[
\sup_{t\ge\tau}
\left|
\mathbb E\!\bigl[e_{B,m}^{D}(\cdot,t)\bigr]
\right|
\le
\varepsilon_{B,m}^{D}(\tau),
\qquad
\sup_{t\ge\tau}
\left|
\mathbb E\!\bigl[e_{B,m}^{H}(\cdot,t)\bigr]
\right|
\le
\varepsilon_{B,m}^{H}(\tau),
\]
with all these functions tending to zero as \(\tau\to\infty\). Define
\[
\varepsilon_{\mathrm{fg}}^{\mathrm{seq}}(\tau)
:=
\sum_{m=1}^{k^*-1}
\Bigl(
\varepsilon_{\mathrm{fg},m}^{\Gamma}(\tau)
+
\varepsilon_{\mathrm{fg},m}^{D}(\tau)
+
\varepsilon_{\mathrm{fg},m}^{H}(\tau)
\Bigr),
\]
\[
\varepsilon_{\mathrm{bg}}^{\mathrm{seq}}(\tau)
:=
\sum_{m=1}^{D-k^*-1}
\Bigl(
\varepsilon_{\mathrm{bg},m}^{\Gamma}(\tau)
+
\varepsilon_{\mathrm{bg},m}^{D}(\tau)
+
\varepsilon_{\mathrm{bg},m}^{H}(\tau)
\Bigr),
\]
with the convention that an empty sum is zero. Then
\[
\sup_{t\ge\tau}
D_{\mathrm{KL}}\!\left(
P_t^{\mathrm{diag}}
\middle\|
\prod_{i=1}^D P_{i,t}
\right)
\le
\varepsilon_\Gamma(\tau)
+
\varepsilon_D(\tau)
+
\varepsilon_H(\tau)
+
\varepsilon_{\mathrm{fg}}^{\mathrm{seq}}(\tau)
+
\varepsilon_{\mathrm{bg}}^{\mathrm{seq}}(\tau).
\]
\end{corollary}

\begin{proof}
The exact block decomposition gives
\[
D_{\mathrm{KL}}\!\left(
P_t^{\mathrm{diag}}
\middle\|
\prod_{i=1}^D P_{i,t}
\right)
=
I_t(u:v)
+
D_{\mathrm{KL}}\!\left(
P_t^{\mathrm{diag}}(u)
\middle\|
\prod_{i\in I^*} P_{i,t}
\right)
+
D_{\mathrm{KL}}\!\left(
P_t^{\mathrm{diag}}(v)
\middle\|
\prod_{j\notin I^*} P_{j,t}
\right).
\]
For any ordered block \(B=(\ell_1,\dots,\ell_r)\), the KL chain rule
gives
\[
D_{\mathrm{KL}}\!\left(
P_{B,t}
\middle\|
\prod_{q=1}^r P_{\ell_q,t}
\right)
=
\sum_{m=1}^{r-1} I_{B,m}(t).
\]
Averaging the assumed prefix decomposition yields
\[
\overline{\Gamma}_{B,m}^{\mathrm{diag}}(t)
=
-I_{B,m}(t)
+
\mathbb E\!\bigl[e_{B,m}^{D}(\cdot,t)\bigr]
+
\mathbb E\!\bigl[e_{B,m}^{H}(\cdot,t)\bigr],
\]
hence
\[
I_{B,m}(t)
\le
\bigl|\overline{\Gamma}_{B,m}^{\mathrm{diag}}(t)\bigr|
+
\left|
\mathbb E\!\bigl[e_{B,m}^{D}(\cdot,t)\bigr]
\right|
+
\left|
\mathbb E\!\bigl[e_{B,m}^{H}(\cdot,t)\bigr]
\right|.
\]
Summing over \(m\) gives the stated bound for each block. Combining
those two blockwise bounds with
Corollary~\ref{cor:conditional_b2b} proves the claim.
\end{proof}

\section{Stable-Hull Transfer and Comparison with HAFF Accessibility}\label{sec:haff}

Let \(\omega(\cdot)=\langle\Omega,\cdot\,\Omega\rangle\) on
\(\mathcal A_U\), and let \(\sigma_s^{\omega,\mathrm{tot}}\) denote the
ambient modular flow. Define
\[
\mathfrak L_{\mathfrak F}^{\mathrm{st}}
:=
\left\{
\begin{aligned}
\mathcal M\subseteq \mathcal A_U:\;&
\mathcal M \text{ is a von Neumann subalgebra},\\
&
\mathcal A_{\mathfrak F}\subseteq \mathcal M,\\
&
\sigma_s^{\omega,\mathrm{tot}}(\mathcal M)=\mathcal M
\quad \forall s\in\mathbb R,\\
&
\mathcal E_t(\mathcal M)\subseteq \mathcal M
\quad \forall t\ge 0
\end{aligned}
\right\}.
\]
This class is nonempty because \(\mathcal A_U\in \mathfrak L_{\mathfrak F}^{\mathrm{st}}\).

\begin{definition}[Stable-hull transfer map]\label{def:stablehull}
The \emph{stable-hull transfer} of the foreground frame
\(\mathfrak F\) is
\[
\Phi_{\mathfrak F}^{\mathrm{st}}(\mathfrak F)
:=
\bigcap_{\mathcal M\in \mathfrak L_{\mathfrak F}^{\mathrm{st}}}\mathcal M.
\]
\end{definition}

\begin{proposition}[Preliminary inclusion of the pointer algebra]\label{prop:prelim_inclusion}
Assume the HAFF standing condition that the accessible algebra is the
fixed-point algebra of the decoherence semigroup,
\[
\mathcal A_c^{\mathrm{HAFF}}
:=
\operatorname{Fix}(\mathcal E)
:=
\{A\in\mathcal A_U:\mathcal E_t(A)=A\ \forall t\ge 0\}.
\]
Let \(\mathfrak F\) be a T-DOME foreground frame, and let
\(O_1,\dots,O_{k^*}\) be exact pointer-aligned observables realizing
\(\iota_{\mathfrak F}\) as in
Proposition~\ref{prop:pointer}. Then
\[
\mathcal A_{\mathfrak F}
=
W^*(1,O_1,\dots,O_{k^*})
\subseteq
\mathcal A_c^{\mathrm{HAFF}}.
\]
\end{proposition}

\begin{proof}
Exact pointer alignment gives
\[
O_i\in \operatorname{Fix}(\mathcal E)=\mathcal A_c^{\mathrm{HAFF}}
\qquad
\text{for each } i.
\]
Since \(\mathcal A_c^{\mathrm{HAFF}}\) is a von Neumann algebra, it
contains the von Neumann algebra generated by \(1,O_1,\dots,O_{k^*}\).
Hence
\[
W^*(1,O_1,\dots,O_{k^*})\subseteq \mathcal A_c^{\mathrm{HAFF}}.
\]
\end{proof}

\begin{proposition}[Minimal stable-hull transfer]\label{prop:stablehull}
Assume the hypotheses of Proposition~\ref{prop:pointer}. Then
\[
\Phi_{\mathfrak F}^{\mathrm{st}}(\mathfrak F)
=
\bigcap_{\mathcal M\in \mathfrak L_{\mathfrak F}^{\mathrm{st}}}\mathcal M
\]
is a well-defined von Neumann subalgebra of \(\mathcal A_U\) satisfying:
\begin{enumerate}
\item \(\mathcal A_{\mathfrak F}\subseteq
  \Phi_{\mathfrak F}^{\mathrm{st}}(\mathfrak F)\);
\item
  \[
  \sigma_s^{\omega,\mathrm{tot}}
  \big(\Phi_{\mathfrak F}^{\mathrm{st}}(\mathfrak F)\big)
  =
  \Phi_{\mathfrak F}^{\mathrm{st}}(\mathfrak F)
  \qquad
  \forall s\in\mathbb R;
  \]
\item
  \[
  \mathcal E_t\big(\Phi_{\mathfrak F}^{\mathrm{st}}(\mathfrak F)\big)
  \subseteq
  \Phi_{\mathfrak F}^{\mathrm{st}}(\mathfrak F)
  \qquad
  \forall t\ge 0;
  \]
\item it is minimal with these properties: if
  \(\mathcal N\subseteq\mathcal A_U\) is any von Neumann subalgebra
  containing \(\mathcal A_{\mathfrak F}\) and stable under both ambient
  dynamics, then
  \[
  \Phi_{\mathfrak F}^{\mathrm{st}}(\mathfrak F)\subseteq \mathcal N.
  \]
\end{enumerate}
\end{proposition}

\begin{proof}
The class \(\mathfrak L_{\mathfrak F}^{\mathrm{st}}\) is nonempty
because \(\mathcal A_U\) belongs to it. An arbitrary intersection of
von Neumann subalgebras is again a von Neumann subalgebra, so the
intersection is well-defined and is itself a von Neumann subalgebra of
\(\mathcal A_U\). Since every element of
\(\mathfrak L_{\mathfrak F}^{\mathrm{st}}\) contains
\(\mathcal A_{\mathfrak F}\), so does the intersection.

If \(A\in\Phi_{\mathfrak F}^{\mathrm{st}}(\mathfrak F)\), then
\(A\in\mathcal M\) for every
\(\mathcal M\in\mathfrak L_{\mathfrak F}^{\mathrm{st}}\). By definition
of the class,
\(\sigma_s^{\omega,\mathrm{tot}}(A)\in\mathcal M\) for every
\(s\in\mathbb R\), and
\(\mathcal E_t(A)\in\mathcal M\) for every \(t\ge 0\). Therefore both
\(\sigma_s^{\omega,\mathrm{tot}}(A)\) and \(\mathcal E_t(A)\) belong to
the intersection. This proves
\(\sigma_s^{\omega,\mathrm{tot}}(\Phi_{\mathfrak F}^{\mathrm{st}}(\mathfrak F))
\subseteq \Phi_{\mathfrak F}^{\mathrm{st}}(\mathfrak F)\). Applying the
same argument with \(-s\) gives the reverse inclusion, hence equality.

Finally, if \(\mathcal N\) is any von Neumann subalgebra with the same
containment and stability properties, then
\(\mathcal N\in\mathfrak L_{\mathfrak F}^{\mathrm{st}}\). Since
\(\Phi_{\mathfrak F}^{\mathrm{st}}(\mathfrak F)\) is the intersection of
all members of that class, it is contained in \(\mathcal N\). This is
exactly the desired minimality.
\end{proof}

\begin{proposition}[Stable-hull transfer lies in the HAFF modular accessible algebra]\label{prop:stablehull_inclusion}
Assume:
\begin{enumerate}
\item the hypotheses of Proposition~\ref{prop:pointer};
\item the exact pointer-invariance condition of
  Proposition~\ref{prop:prelim_inclusion};
\item there exists a HAFF modular accessible algebra
  \(\mathcal A_c^{\mathrm{HAFF}}\subseteq \mathcal A_U\) which is
  invariant under the ambient modular flow and forward-stable under the
  decoherence semigroup in the sense of HAFF Paper~D~\cite{HAFFD}.
\end{enumerate}
Then
\[
\Phi_{\mathfrak F}^{\mathrm{st}}(\mathfrak F)\subseteq
\mathcal A_c^{\mathrm{HAFF}}.
\]
\end{proposition}

\begin{proof}
By Proposition~\ref{prop:prelim_inclusion},
\(\mathcal A_{\mathfrak F}\subseteq \mathcal A_c^{\mathrm{HAFF}}\).
Assumption~(3) says that \(\mathcal A_c^{\mathrm{HAFF}}\) is a von
Neumann subalgebra, invariant under the ambient modular flow, and
forward-stable under the decoherence semigroup. Hence
\(\mathcal A_c^{\mathrm{HAFF}}\in \mathfrak L_{\mathfrak F}^{\mathrm{st}}\).
Since \(\Phi_{\mathfrak F}^{\mathrm{st}}(\mathfrak F)\) is the
intersection of all members of that class, it is contained in
\(\mathcal A_c^{\mathrm{HAFF}}\).
\end{proof}

\begin{corollary}[Partial Lemma~C]\label{cor:partialC}
Assume the hypotheses of Proposition~\ref{prop:stablehull}. Then
\(\Phi_{\mathfrak F}^{\mathrm{st}}(\mathfrak F)\) is a von Neumann
subalgebra of \(\mathcal A_U\), and the restricted vector state
\[
\omega_{\mathfrak F}^{\mathrm{st}}
:=
\langle\Omega,\cdot\,\Omega\rangle|_{\Phi_{\mathfrak F}^{\mathrm{st}}(\mathfrak F)}
\]
is faithful. If, in addition, the hypotheses of
Proposition~\ref{prop:stablehull_inclusion} hold, then
\[
\Phi_{\mathfrak F}^{\mathrm{st}}(\mathfrak F)\subseteq
\mathcal A_c^{\mathrm{HAFF}}.
\]
\end{corollary}

\begin{proof}
The subalgebra statement is exactly
Proposition~\ref{prop:stablehull}. Because \(\Omega\) is separating for
the ambient algebra \(\mathcal A_U\), it is also separating for every
von Neumann subalgebra of \(\mathcal A_U\); hence the restricted vector
state is faithful on
\(\Phi_{\mathfrak F}^{\mathrm{st}}(\mathfrak F)\). The inclusion claim
is Proposition~\ref{prop:stablehull_inclusion}.
\end{proof}

\paragraph{Anchored shell existence around the stable hull.}
The stable-hull output already anchors a Foundation~I shell problem.
Write
\[
\mathcal A_*:=\Phi_{\mathfrak F}^{\mathrm{st}}(\mathfrak F).
\]
For \(X\in\operatorname{Ball}(\mathcal A_U)\) and \(q\ge 0\), define
the \(\mathcal A_*\)-relative defect
\[
\delta_q^{\mathcal A_*,\Omega}(X)
:=
\operatorname{dist}_{\Omega}^{\sharp}
\bigl(\mathcal E_q(X),\,\operatorname{Ball}(\mathcal A_*)\bigr),
\qquad
\|Y\|_\Omega^\sharp
:=
\bigl(\|Y\Omega\|^2+\|Y^*\Omega\|^2\bigr)^{1/2}.
\]
For \(\varepsilon\ge 0\) and \(\tau>0\), let
\[
\mathfrak R_{\varepsilon,\tau}^{\mathcal A_*,\Omega}
:=
\left\{
\mathcal M\ \middle|\
\begin{array}{l}
\mathcal A_*\subseteq \mathcal M\subseteq \mathcal A_U,\ \mathcal M
\text{ is a unital vN subalgebra},\\[1mm]
\delta_q^{\mathcal A_*,\Omega}(X)\le \varepsilon,\ \forall
X\in\operatorname{Ball}(\mathcal M),\
\forall q\in\mathbb Q\cap[0,\tau]
\end{array}
\right\}.
\]

\begin{proposition}[Anchored relative-shell existence around the stable hull]\label{prop:anchored_shell}
Assume the hypotheses of Corollary~\ref{cor:partialC}. Then for every
\(\varepsilon\ge 0\) and \(\tau>0\):
\begin{enumerate}
\item the class
  \(\mathfrak R_{\varepsilon,\tau}^{\mathcal A_*,\Omega}\) is
  nonempty;
\item every chain in
  \(\mathfrak R_{\varepsilon,\tau}^{\mathcal A_*,\Omega}\) has an
  upper bound;
\item therefore maximal elements exist.
\end{enumerate}
\end{proposition}

\begin{proof}
By Corollary~\ref{cor:partialC}, \(\mathcal A_*\) is a von Neumann
subalgebra of \(\mathcal A_U\) and is forward-stable under the
decoherence semigroup. Hence \(\mathcal A_*\) itself belongs to
\(\mathfrak R_{\varepsilon,\tau}^{\mathcal A_*,\Omega}\): for
\(X\in\operatorname{Ball}(\mathcal A_*)\), forward stability gives
\(\mathcal E_q(X)\in\mathcal A_*\), and contractivity gives
\(\mathcal E_q(X)\in\operatorname{Ball}(\mathcal A_*)\), so
\(\delta_q^{\mathcal A_*,\Omega}(X)=0\). Thus the class is nonempty.

Let \(\{\mathcal M_i\}_{i\in I}\) be a chain in
\(\mathfrak R_{\varepsilon,\tau}^{\mathcal A_*,\Omega}\), and set
\[
\mathcal M:=W^*\!\Big(\bigcup_{i\in I}\mathcal M_i\Big).
\]
Then \(\mathcal A_*\subseteq \mathcal M\subseteq \mathcal A_U\), and
\(\mathcal M\) is a unital von Neumann subalgebra. By Kaplansky
density, every \(X\in\operatorname{Ball}(\mathcal M)\) is a bounded
strong* limit of a net from the unit ball of the algebraic union
\(\bigcup_i\mathcal M_i\). Since the chain is totally ordered, each
approximant lies in some admissible member of the chain, so its
\(\mathcal A_*\)-relative defect is bounded by \(\varepsilon\) for all
\(q\in\mathbb Q\cap[0,\tau]\). Passing to the limit in the
\(\Omega^\sharp\)-seminorm yields the same bound for \(X\), hence
\(\mathcal M\in\mathfrak R_{\varepsilon,\tau}^{\mathcal A_*,\Omega}\).
Therefore every chain has an upper bound. Zorn's lemma now gives a
maximal element.
\end{proof}

\begin{remark}[Status of anchored shells]\label{rem:anchored_shell_status}
Proposition~\ref{prop:anchored_shell} does not yet construct the full
shell family required on the Paper~I side. It closes only the anchored
existence problem around the theorem-level Foundation~I output
\(\mathcal A_*\). What remains open is the upgrade to a canonical,
nontrivial, and calibrated shell family with the correct modular
transport law.
\end{remark}

\begin{remark}[Why equality is not the right first target]
In general, \(\mathcal A_c^{\mathrm{HAFF}}\) may be much larger than the
minimal modular-decoherence stable hull of finitely many foreground
pointer observables. The honest first target is therefore the inclusion
\[
\Phi_{\mathfrak F}^{\mathrm{st}}(\mathfrak F)\subseteq
\mathcal A_c^{\mathrm{HAFF}},
\]
not the equality
\[
\Phi_{\mathfrak F}^{\mathrm{st}}(\mathfrak F)\stackrel{?}{=}
\mathcal A_c^{\mathrm{HAFF}}.
\]
The equality problem should remain open until one proves that the
foreground realization together with the stable-hull transfer selects
the same algebra as the HAFF maximality construction.
\end{remark}

\begin{remark}[Relation to Paper~III]
One should not identify the foreground complexity \(k^*\) directly with
the statistical complexity parameter \(n\) of Paper~III~\cite{PaperIII}.
The safer statement is only that
\(\Phi_{\mathfrak F}^{\mathrm{st}}(\mathfrak F)\) provides a natural
algebraic candidate for a later bridge to the finite-observer sector
studied there.
\end{remark}

\section{What Remains Open}\label{sec:open}

The bottlenecks are now sharper than they were before this note.

\begin{enumerate}
\item \textbf{Persistent projection selection.} Lemma~\ref{lem:A} is
  only conditional. The constructive question of which projection
  emerges from the source datum remains open.
\item \textbf{Foreground realization and model-independent PS2.} The
  stable-hull route closes the transfer-map well-definedness problem
  once \(\mathcal A_{\mathfrak F}\) exists. Within the restricted
  product-classical and finite-dimensional faithful pointer regime, the
  PS2 package is now theorem-backed:
  Lemma~\ref{lem:product_classical} closes the structural reduction to
  fragment-local omission losses, Lemma~\ref{lem:fragment_orth} closes
  orthogonality, and Corollary~\ref{cor:ps2_closure} packages the
  restricted selection theorem. Lemma~\ref{lem:diag_off_split} now
  fixes the bookkeeping split of the survival deficit into a
  pointer-diagonal term and an off-sector remainder. In addition,
  Proposition~\ref{prop:conditional_b2a} and
  Corollary~\ref{cor:conditional_b2b} now show that the
  foreground-background part of the model-independent programme closes
  conditionally once four tail inputs are granted: approximate tail
  optimality, exchange-cover plus persistence, tail-passive background
  survival, and an exchange-efficiency gap, together with the
  block-locality defects for the KL and control terms. Corollary~
  \ref{cor:conditional_full_prod} then reduces the remaining diagonal
  product-defect problem to finitely many ordered within-block prefix
  couplings. The remaining open issue is therefore to derive those tail
  hypotheses together with the prefix-level coupling and locality
  bounds model-independently from the broken-phase output and the full
  T-DOME survival identity rather than assume them. Beyond the
  restricted regime, the unresolved analytic obstacles are explicit:
  cross-channel correlations need not decouple, pointer spectra need
  not remain finite or discrete, and the Hilbert--Schmidt control used
  here is not available in general Type~III infinite-dimensional
  settings, where a predual reformulation may be necessary. The
  pointer-sector existence input PS1 is treated here only as an
  exact-regime idealization; what remains genuinely open is the
  model-independent alignment step PS2 and the induced ordered-selection
  step PS3.
\item \textbf{Comparison with HAFF maximality.} The present paper proves
  inclusion only. Equality with the full HAFF accessible algebra, or any
  canonical equivalence theorem, remains a separate strong-form problem.
\item \textbf{Bridge to Paper~I admissibility.} To feed directly into
  Paper~I, one must still verify that the algebra produced by the lower
  bridge satisfies the additional admissibility hypotheses used there,
  beyond the already-closed faithfulness condition. More precisely, the
  remaining interface is to derive the modular-decoherence scaling axiom
  \((C)\) together with a Foundation~I shell family
  \(\{\mathcal A_*(\tau)\}_{\tau>0}\) around the stable-hull output
  \(\mathcal A_*:=\Phi_{\mathfrak F}^{\mathrm{st}}(\mathfrak F)\), and
  one-scale cyclicity of the corresponding shell-annulus algebra
  \[
  \mathcal D_{\tau}:=\mathcal A_*(\tau)\cap \mathcal A_*'.
  \]
  In Paper~I language these are exactly the remaining inputs
  \((C)+(Amb)+(CN_\exists)\). The shell-family side has now itself been
  partially internalized: Proposition~\ref{prop:anchored_shell} closes
  anchored relative-shell existence around \(\mathcal A_*\). What
  remains open is the canonical and calibrated nontrivial shell family.
  The naive scalar-threshold route now has an exact blocker: fixed
  thresholds are compatible with modular transport but not with exact
  core calibration, vanishing thresholds give exact core calibration but
  not automatic modular transport, and the only scalar profile with
  both properties is the trivial zero-threshold profile, which
  collapses the shells back to \(\mathcal A_*\). More sharply, even the
  fixed-threshold branch now reduces to finding a scale-invariant
  nontriviality window between core collapse and ambient collapse,
  together with uniqueness of maximal shells inside that window; exact
  calibration is a separate issue beyond that.
\end{enumerate}

These open problems are real, but they are no longer vague. Foundation~I
is now separated into one conditional source-selection problem,
one theorem-level thermodynamic packaging result, one theorem-level
stable-hull transfer construction, and one remaining
foreground-realization/comparison problem.

\section{Conclusion}

The purpose of this paper has been calibration, not inflation.
Foundation~I is neither absent nor complete. It is partially available
in theorem-level form, partially available as an explicit construction
problem, and partially still conditional.

The main positive result is that the T-DOME I--II chain can already be
written as a formal Foundation~I proposition, and that the first
restricted PS2 bottleneck can now be closed at theorem level in a first
restricted regime. In the product-classical, decoherence-dominated
pointer regime, pointer coarse-graining yields fragment-local omission
losses, independent fragment records force orthogonality, semigroup
decay plus explicit Hilbert--Schmidt continuity constants bound the
coherence remainder, and a sufficiently large weight gap forces exact
top-\(k^*\) pointer selection. On the algebraic side, the first
operator-algebraic realization of the foreground sector can now be
extended to a well-defined minimal modular-decoherence stable hull,
yielding a faithful von Neumann subalgebra together with the inclusion
\[
\Phi_{\mathfrak F}^{\mathrm{st}}(\mathfrak F)\subseteq
\mathcal A_c^{\mathrm{HAFF}}.
\]
At the level of future PS2 work, the paper now also fixes the
bookkeeping split of the survival deficit into a pointer-diagonal term
and an off-sector remainder, and it closes the first
foreground-background decoupling step conditionally: if the tail
distortion localizes, the exchange-weighted defect persists, the
background diagonal survival becomes tail-passive, the exchange
efficiency ratio stays below the squared-loss threshold, and the
remaining KL/control locality defects are small, then the broken phase
forces tail-small diagonal foreground-background coupling and hence
tail-small foreground-background mutual information. Corollary~
\ref{cor:conditional_full_prod} then shows that the unresolved
product-defect part is no longer a monolithic statement: it is a finite
family of ordered within-block prefix decoupling problems.
The remaining model-independent problem is therefore no longer vague: it
is to derive those tail hypotheses and the ordered prefix-level
coupling/locality bounds, and thereby to control the full diagonal
product defect as a theorem-level output of the broken phase.
On the Paper~I side, the admissibility interface is now sharpened in the
same way: faithfulness is already closed, ambient scaling covariance
would transfer the dynamical axiom \((C)\), and the remaining regularity
burden is exactly the existence of a Foundation~I shell family and one
cyclic shell-annulus scale. Even the shell-family side is no longer a
black box: anchored relative shells around \(\mathcal A_*\) exist, but
their upgrade to a canonical nontrivial shell family currently stops at
the scalar-threshold trilemma between transport, exact core
calibration, and nontriviality. Within the fixed-threshold route, this
has already reduced further to a scale-invariant nontriviality window
and uniqueness problem, plus a separate calibration theorem.
The main negative result is equally important: nothing in the present
note licenses the stronger claim that the source datum has already been
shown to generate the full accessible algebra uniquely, nor that PS2 has
been derived model-independently from the broken-phase output. Those
remain the unfinished parts of the unified-source programme.

In that sense, the current status is simple. The thread can be strung,
but only honestly: conditional at the source end, inclusion-first at the
algebra end, and theorem-backed only where the intermediate steps have
actually been closed.

\appendix

\section{External Genesis Inputs (Self-Contained Statements)}\label{app:premises}

The following list records the external Genesis-side inputs used in the
paper in a self-contained form. The point is not to reproduce their full
proofs, but to state exactly what is imported from HAFF and T-DOME.

\begin{enumerate}
\item \textbf{HAFF ambient datum.}
The ambient setting is a Type~III\(_1\) von Neumann algebra in standard
form \((\mathcal H_U,\mathcal A_U,\Omega)\), with \(\Omega\) cyclic and
separating for \(\mathcal A_U\), together with a
\FIAppendixContinuityPhrase\ semigroup of normal completely positive
unital maps \(\{\mathcal E_t\}_{t\ge 0}\) on \(\mathcal A_U\). This is
the source datum formalized in Section~\ref{sec:ledger}.

\item \textbf{T-DOME~I projected record dynamics.}
Given a persistent Nakajima--Zwanzig projection
\(\mathcal P:\mathcal A_U\to\mathcal A_U\), the projected predual
dynamics is analyzed in a non-Markovian kernel form
\[
\partial_t \omega_{\mathcal P,t}
=
\int_0^t \mathcal K_*(t,s)\,\omega_{\mathcal P,s}\,ds
\;+\; \text{inhomogeneous term},
\]
or equivalently in the dual memory-kernel formulation used in
Section~\ref{sec:ledger}. Foundation~I uses only the existence of such a
kernel and a persistent record process, not an explicit closed-form
solution.

\item \textbf{T-DOME~I survival identity and distortion functional.}
The survival bookkeeping imported into the PS2 analysis is the identity
\[
\beta\,\mathcal S
=
-\Delta I(S{:}E)
-\Delta D_{\mathrm{KL}}
-\beta\,\Delta\langle H_{\mathrm{ctrl}}\rangle,
\]
with frame comparisons measured in the squared-loss distortion
\[
D(\mathcal F_I)
:=
\mathbb E_\xi\!\left[
\bigl(
\mathcal S_{\mathrm{full}}(\xi)
-\mathcal S_{\mathcal F_I}(\xi)
\bigr)^2
\right].
\]
The diagonal/off-sector split and all restricted PS2 estimates are built
from this identity rather than assumed independently.

\item \textbf{T-DOME~II symmetry-broken compression.}
If a persistent record process has effective dimension \(D\) and
per-component entropy rate \(h_\mu\), then the symmetric full-rank
processing requirement is
\[
R_{\mathrm{full}}=h_\mu D.
\]
When
\[
R_{\mathrm{full}}
>
\dot{\mathcal C}_{\mathrm{budget}},
\]
the full-rank symmetric strategy is infeasible. Under the standing
non-degeneracy hypotheses on the survival-distortion landscape, the
survival-optimal admissible rank is
\[
k^*
=
\Big\lfloor
\frac{\dot{\mathcal C}_{\mathrm{budget}}}{h_\mu}
\Big\rfloor
< D,
\]
and a symmetry-broken foreground frame
\(\mathfrak F=\{f_1,\dots,f_{k^*}\}\) emerges, ordered by survival
relevance.

\item \textbf{Decoherence-side pointer sector (exact-regime idealization).}
In the exact dephasing-dominated regime, robust classical records are
carried by a commuting pointer sector. Foundation~I uses this as the
PS1 input: a commuting family of exact fixed-point projections for
\(\{\mathcal E_t\}\) generating a classical pointer algebra inside
\(\mathcal A_U\). This is an idealization of the approximate
einselection and redundant-record structure described in the quantum
Darwinism literature~\cite{Zurek2003,Zurek2009}, not a consequence of
generic decoherence alone.

\item \textbf{HAFF accessible algebra as stable exact sector.}
In the exact fixed-point version used for the inclusion comparison, the
HAFF accessible algebra is an ambient von Neumann subalgebra
\(\mathcal A_c^{\mathrm{HAFF}}\subseteq \mathcal A_U\) that is
invariant under the ambient modular flow and forward-stable under the
decoherence semigroup. In the strong exact form,
\(\mathcal A_c^{\mathrm{HAFF}}=\operatorname{Fix}(\mathcal E)\).

\item \textbf{T-DOME~III / Paper~III calibration distinction.}
Foundation~I uses T-DOME~III and Paper~III only as contextual support
for what is \emph{not} claimed: the foreground complexity \(k^*\) is not
identified with Paper~III's statistical complexity parameter \(n\), and
the present paper does not attempt a direct derivation of Paper~III's
self-description bounds from the source-side chain.
\end{enumerate}

\section{Conditional Input Ledger}\label{app:conditional_ledger}

For ease of reference, the remaining conditional interfaces are listed
here in one place.

\begin{enumerate}
\item \textbf{Source selection.}
Lemma~\ref{lem:A} is conditional on the existence of a persistent
Nakajima--Zwanzig projection producing the projected record kernel.
This is the observer-emergence bottleneck.

\item \textbf{Pointer alignment package.}
The foreground realization problem separates into:
\begin{enumerate}
\item[(PS1)] existence of a decoherence-selected exact fixed-point
  pointer sector (exact dephasing-dominated regime idealization, not a
  consequence of generic decoherence);
\item[(PS2)] identification of the T-DOME sufficient-statistic
  coordinates with the pointer-sector coordinates;
\item[(PS3)] ordered pointer-channel selection by survival relevance.
\end{enumerate}
PS1 is an exact-regime idealization of the decoherence-side input; the
genuine model-independent open problem is PS2, with PS3 downstream
from it.

\item \textbf{Fixed-point sector structural inputs.}
Several results in the paper use the following structural properties of
the decoherence semigroup, which hold in the exact dephasing-dominated
regime but are not automatic for a general normal UCP semigroup on a
Type~III\(_1\) factor:
\begin{enumerate}
\item \(\operatorname{Fix}(\mathcal E)\) is a von Neumann subalgebra of
  \(\mathcal A_U\);
\item there exists a faithful normal conditional expectation
  \(\mathbb E_{\mathrm{ptr}}:\mathcal A_U\to
  \mathcal A_{\mathrm{ptr}}\) preserving \(\omega\);
\item the pointer algebra \(\mathcal A_{\mathrm{ptr}}\) is contained in
  \(\operatorname{Fix}(\mathcal E)\).
\end{enumerate}
These are standing inputs of the exact-regime framework used throughout
Foundation~I, not consequences of the source datum
alone~\cite{Takesaki1972,FrigerioVerri1982}.

\item \textbf{Restricted PS2 regime assumptions.}
The theorem-level restricted PS2 closure is proved only in the
product-classical, finite-dimensional faithful pointer regime stated in
Section~\ref{sec:ps2model}. In particular, the restricted closure does
not yet treat general cross-channel correlations, continuous pointer
spectra, or the genuinely infinite-dimensional Type~III estimate
problem.

\item \textbf{Conditional foreground-background decoupling inputs.}
Proposition~\ref{prop:conditional_b2a} requires the following named
inputs on the tail regime: approximate tail optimality, exchange-cover
plus persistence, tail-passive background survival, and an
exchange-efficiency gap, together with the KL/control locality defects
used in Corollary~\ref{cor:conditional_b2b}.

\item \textbf{Conditional full product-defect inputs.}
Corollary~\ref{cor:conditional_full_prod} adds the ordered within-block
prefix-chain analogues of the same decoupling/locality inputs. The full
diagonal product-defect problem is therefore reduced to finitely many
named prefix couplings rather than left as one undifferentiated gap.

\item \textbf{Paper~I admissibility inputs.}
The lower output feeds into Paper~I only conditionally on three further
interfaces: ambient modular-decoherence scaling covariance, a
Foundation~I shell family around the stable hull, and one-scale
shell-annulus cyclicity for that family. Proposition~
\ref{prop:anchored_shell} closes only the anchored existence part of
this shell problem.
\end{enumerate}

\newpage
\addcontentsline{toc}{section}{References}

\begin{thebibliography}{99}

\bibitem{PaperI}
S.~Liu,
\emph{The Inseparability of Geometry and Modular Flow: A Structural Inseparability Theorem and Its Implications for the Hard Problem},
Consciousness Trilogy, Paper~I (2026),
Zenodo, DOI: 10.5281/zenodo.19607387.

\bibitem{PaperIII}
S.~Liu,
\emph{The Structural Horizon of Self-Description: Fisher Bounds, Small-Ball Remainders, and the Quantitative Explanatory Gap},
Consciousness Trilogy, Paper~III (2026),
Zenodo, DOI: 10.5281/zenodo.19607724.

\bibitem{HAFFD}
S.~Liu,
\emph{Gravitational Phenomena as Emergent Properties of Observable Algebra Selection: A Structural Analysis},
Zenodo (2026), DOI: 10.5281/zenodo.18388881.

\bibitem{TDOMEI}
S.~Liu,
\emph{Non-Markovian Memory and the Thermodynamic Necessity of Temporal Accumulation},
Zenodo (2026), DOI: 10.5281/zenodo.18574342.

\bibitem{TDOMEII}
S.~Liu,
\emph{Spontaneous Symmetry Breaking of Reference Frames as a Computational Cost Minimization Strategy},
Zenodo (2026), DOI: 10.5281/zenodo.18579703.

\bibitem{TDOMEIII}
S.~Liu,
\emph{Fisher Information Geometry and the Thermodynamic Cost of Self-Referential Calibration},
Zenodo (2026), DOI: 10.5281/zenodo.18591771.

\bibitem{BreuerPetruccione2002}
H.-P.~Breuer and F.~Petruccione,
\emph{The Theory of Open Quantum Systems},
Oxford University Press (2002).

\bibitem{Nakajima1958}
S.~Nakajima,
\emph{On quantum theory of transport phenomena},
Prog.\ Theor.\ Phys.\ \textbf{20}, 948 (1958).

\bibitem{Zurek2003}
W.~H.~Zurek,
\emph{Decoherence, einselection, and the quantum origins of the classical},
Rev.\ Mod.\ Phys.\ \textbf{75}, 715 (2003).

\bibitem{Zurek2009}
W.~H.~Zurek,
\emph{Quantum Darwinism},
Nature Physics \textbf{5}, 181 (2009).

\bibitem{Zwanzig1960}
R.~Zwanzig,
\emph{Ensemble method in the theory of irreversibility},
J.\ Chem.\ Phys.\ \textbf{33}, 1338 (1960).

\bibitem{Takesaki1972}
M.~Takesaki,
\emph{Conditional expectations in von Neumann algebras},
J.\ Funct.\ Anal.\ \textbf{9}, 306--321 (1972).

\bibitem{FrigerioVerri1982}
A.~Frigerio and M.~Verri,
\emph{Long-time asymptotic properties of dynamical semigroups on $W^*$-algebras},
Math.\ Z.\ \textbf{180}, 275--286 (1982).

\end{thebibliography}

\end{document}
